\documentclass[11pt,a4paper]{article}
\usepackage[T1]{fontenc}
\usepackage{lmodern,microtype,amsmath,amssymb,amsthm,mathtools}
\usepackage[margin=24mm,headheight=15pt]{geometry}
\usepackage{booktabs,array,tabularx,xcolor,fancyhdr,titlesec}
\usepackage[colorlinks=true,linkcolor=teal!65!black,citecolor=teal!65!black,urlcolor=teal!65!black]{hyperref}
\definecolor{ink}{HTML}{15364A}
\definecolor{accent}{HTML}{007C78}
\titleformat{\section}{\Large\sffamily\bfseries\color{ink}}{\thesection}{.55em}{}
\titleformat{\subsection}{\large\sffamily\bfseries\color{accent}}{\thesubsection}{.55em}{}
\newtheorem{theorem}{Theorem}[section]
\newtheorem{lemma}[theorem]{Lemma}
\newtheorem{proposition}[theorem]{Proposition}
\newtheorem{corollary}[theorem]{Corollary}
\theoremstyle{remark}\newtheorem{remark}[theorem]{Remark}
\newtheorem{example}[theorem]{Example}
\DeclareMathOperator{\conv}{conv}
\DeclareMathOperator{\vol}{vol}
\DeclareMathOperator{\inter}{int}
\DeclareMathOperator{\relint}{relint}
\newcommand{\R}{\mathbb R}\newcommand{\Z}{\mathbb Z}
\pagestyle{fancy}\fancyhf{}
\fancyhead[L]{\small\sffamily Arithmetic gaps and covering products}
\fancyhead[R]{\small\sffamily Proof manuscript}
\fancyfoot[L]{\footnotesize\sffamily Prepared for Keio | 20 November 2026}
\fancyfoot[R]{\small\thepage}
\setlength{\parindent}{1.1em}\setlength{\parskip}{3pt}
\setlength{\emergencystretch}{2em}\raggedbottom
\hypersetup{pdftitle={Arithmetic Gaps, Exact Minima, and Covering Products},pdfauthor={Research manuscript prepared for Keio}}


\begin{document}
\thispagestyle{empty}
\noindent{\sffamily\small\bfseries\color{accent}
GEOMETRY OF NUMBERS | NEW TRANSFERS FROM VMRC [200]}\par\bigskip
\noindent{\sffamily\bfseries\fontsize{25}{29}\selectfont\color{ink}\raggedright
Arithmetic Gaps, Exact Minima,\\
and Covering Products\par}
\medskip
\noindent{\large\sffamily Four additional results from the\\
Verifiable Mathematical Reasoning Corpus}\par
\bigskip
\noindent Research manuscript prepared for Keio\hfill 20 November 2026
\bigskip
\begin{abstract}
We transfer several elementary certificate mechanisms from the corpus
to problems left open by the preceding lattice-covering notes.
A residue deficit at the jump events of a fractional-part flow gives
an arithmetic improvement involving
$\sum_j\gcd(Q,q_j)$ for a positive simplex relation $q$.
It proves the numerical interior-point conjecture for all non-hollow
lattice simplices of prime normalized volume, with no collinearity
assumption, and gives the exact radius for every relation
$(1^k,a^\ell)$. These include infinite families in the previously
untreated two-interior-point regime.
For arbitrary real $0<a\le b$, $c\ge2b$, and every $n\ge2$, we
determine $\mu_2(S_n(a,b,c,\ldots,c))$ by an explicit maximum of
$n-1$ expressions; a circular pigeonhole argument makes this finite
formula sharp in the given ambient dimension.
We prove supermultiplicativity of the covering product under arbitrary
convex coordinatewise-monotone combinations of factor gauges, extending
both free sums and Cartesian products.
Finally, every extremal non-hollow $d$-polytope has at most $d!$
last-covered classes in its lattice torus, with equality precisely
for the standard terminal simplex. Nonsimplicial polytopes have the
sharp stronger bound $(d-1)!$. A recursive finite certificate decides
extremality for every non-hollow lattice polytope.
Complete proofs and explicit remaining limitations are supplied.
\end{abstract}
\vfill
\noindent\begin{minipage}{\linewidth}\small
\textbf{Status and scope.} This is a new research note, separate from
the preceding 23-, 27-, and 15-page notes. The unrestricted numerical
interior-point conjecture and covering-product conjecture remain
unresolved here. The finite equality test is distinguished from a
geometric classification, and the three-weight formula is stated with
its actual parameter restriction. Proofs, independent mathematical
checks, and exact arithmetic examples support the statements below;
external specialist review and historical priority are separate matters.
\end{minipage}
\clearpage
\tableofcontents
\clearpage

\section{Results, conventions, and the transfer map}

For a full-dimensional compact convex body $K\subset\R^d$ and a
full-rank lattice $\Lambda$, let
\[
 \mu(K,\Lambda)=\inf\{t\ge0:tK+\Lambda=\R^d\}.
\]
The $i$th covering minimum is the least dilation whose lattice
translates meet every affine subspace of codimension $i$. We use
its standard rational-quotient characterization \cite{CSS,Krivokuca}:
\[
 \mu_i(K,\Lambda)=\sup_A\mu(AK,A\Lambda),
\]
where $A$ has rank $i$ and lattice-rational kernel. All section and
quotient lattices are the actual induced lattices. A non-hollow
lattice polytope has at least one interior lattice point.

Write $h=|\operatorname{int}P\cap\Lambda|$. The numerical part of
Codenotti--Santos--Schymura Conjecture D is \cite{CSS}
\begin{equation}\label{intro:D}
 \mu(P,\Lambda)\le\frac{d-1}{2}+\frac1{h+1}.
\end{equation}
For a general convex body, Conjecture F is \cite{GMS,CSS}
\begin{equation}\label{intro:F}
 \mathcal F_d(K,\Lambda):=
 \frac{\operatorname{vol}_d(K)}{\det\Lambda}
       \prod_{i=1}^{d}\mu_i(K,\Lambda)
 \ge\frac{d+1}{2^d}.
\end{equation}
Here $\operatorname{vol}_d/\det\Lambda$ is the ordinary
lattice-normalized Euclidean volume. By contrast, the
\emph{normalized volume of a lattice simplex} in the prime-volume
result is $d!\operatorname{vol}_d/\det\Lambda$, so a unimodular
simplex has normalized volume one.

For positive real weights define
\[
 S_n(\omega_0,\ldots,\omega_n)=
 \operatorname{conv}(-\omega_0\mathbf1_n,\omega_1e_1,\ldots,\omega_ne_n),
 \qquad T_n=S_n(1,\ldots,1).
\]

\begin{center}\small
\begin{tabularx}{\linewidth}{@{}>{\raggedright\arraybackslash}p{.15\linewidth}
 >{\raggedright\arraybackslash}X >{\raggedright\arraybackslash}X@{}}
\toprule
Target & Result established in this note & Remaining limitation\\
\midrule
D & Prime-volume simplices; coprime relation numerators; exact radii for
$(1^k,a^\ell)$ & Arbitrary polytopes and all of the two-point regime\\
Weighted $\mu_2$ & Every $S_n(a,b,c,\ldots,c)$ with $0<a\le b$ and $c\ge2b$ &
Arbitrary independent weights, $b<c<2b$, and higher indices\\
F & Arbitrary monotone gauge combinations; every combination of planar
bodies and intervals & General bodies without this structure\\
Equality in A & Sharp deep-hole counts and a finite test for every polytope &
A geometric list of all extremal polytopes\\
\bottomrule
\end{tabularx}
\end{center}

\subsection{Which corpus mechanisms are actually transferred}

We inspected the corpus at revision
\texttt{ac1492a50784be919fe0b5e8f19f96f8d221a22e}, also confirmed
as its current main revision during this work. The transfers use
specific arguments rather than assuming unrelated corpus claims.

\begin{center}\small
\begin{tabularx}{\linewidth}{@{}>{\raggedright\arraybackslash}p{.19\linewidth}
 >{\raggedright\arraybackslash}X >{\raggedright\arraybackslash}X@{}}
\toprule
Corpus source & Mechanism read & New use here\\
\midrule
\texttt{[030]} & Endpoint-aware finite sawtooth distribution &
Arithmetic deficits at the events of [200]\\
\texttt{[029]}, \texttt{[215]} & Torus pigeonhole returns; weighted partition of phases &
Exact two-event placement and the $n-1$ multiplicity bound\\
\texttt{[185]}, \texttt{[123]} & Nonnegative local packing certificates; beta-integral
and gamma-factor comparisons &
Permutation-simplex packing and gauge-product inequalities\\
\texttt{[041]} & Conditional finite states composed through fibers &
Recursive lifting of extremal deep-hole states\\
\bottomrule
\end{tabularx}
\end{center}

The finite ingredients needed in the new geometry are proved explicitly
below. The cited corpus mechanisms are documented in
\cite{Corpus030,Corpus029,Corpus215,Corpus185,Corpus123,Corpus041};
their other research claims are not premises. The basic flow and
non-hollow bound of \cite{VMRC200}, the earlier simplex congruence
criterion, and the earlier one-small-column lemma are reproduced
as inputs, not counted as additional results.


\section{Covering tools used in the proofs}

\begin{lemma}[Weighted fractional-part descent]\label{base:flow}
Let $q_j$ be positive integers, $c_j>0$, and $t_j\in\R$. Put
$C=\sum_jq_jc_j$ and $F(s)=\sum_jc_j\{t_j+sq_j\}$. Then
\[
 \min_s F(s)\le
 \frac12\sum_jc_j-\frac{\sum_jq_jc_j^2}{2C}.
\]
\end{lemma}
\begin{proof}
Between its finitely many events on a period, $F$ has slope $C$;
its mean is $\frac12\sum_jc_j$. At an event $e$ let $J_e$ be
the total downward jump and let $F(e)$ be the post-jump value.
Integrating the continuous derivative of $F^2$, and adding its
jumps, gives
\[
 \frac1C\sum_eJ_eF(e)
 =\frac12\sum_jc_j-\frac1{2C}\sum_eJ_e^2.
\]
Here $\sum_eJ_e=C$ and
$\sum_eJ_e^2\ge\sum_jq_jc_j^2$, since simultaneous positive
jumps only increase their squared total. The minimum is at most
the weighted average on the left.
\end{proof}

\begin{proposition}[Primitive-direction simplex bound]\label{base:simplex}
Let $S=\conv(v_0,\ldots,v_r)$ contain zero relatively internally,
write $v_j=\lambda_j u_j$ with $u_j$ primitive in its ambient lattice
and $\lambda_j>0$, and let $\alpha_j$ be the barycentric coordinates
of zero. Then
\[
 \mu(S)\le\frac12\sum_{j=0}^{r}\frac{1-\alpha_j}{\lambda_j}.
\]
\end{proposition}
\begin{proof}
The integer directions have a positive integer relation $q$.
Put $c_j=1/\lambda_j$ and $C=\sum_jq_jc_j$; then
$\alpha_j=q_jc_j/C$. Given $x=\sum_j t_j u_j$, subtract the lattice
vector $\sum_j\lfloor t_j+sq_j\rfloor u_j$. The remainder is
$\sum_jc_j\{t_j+sq_j\}v_j$, a nonnegative combination of the vertices.
Lemma~\ref{base:flow} bounds its coefficient sum by
$\frac12\sum_j(1-\alpha_j)c_j$. This proves the covering bound.
Only membership of the directions in the lattice is required, not
generation of that lattice.
\end{proof}

\begin{lemma}[Section and quotient]\label{base:section}
If $0\in\operatorname{int}K$ and $L$ is a lattice-rational linear
subspace with quotient map $\pi$, then
\[
 \mu(K,\Lambda)\le
 \mu(K\cap L,\Lambda\cap L)+\mu(\pi K,\pi\Lambda).
\]
\end{lemma}
\begin{proof}
Choose $a,b$ strictly above the two radii. For any $x$, use the
quotient covering to find $z\in\Lambda$ and $y\in bK$ with
$\pi(x-z)=\pi y$. Cover $x-z-y\in L$ by a translate of
$a(K\cap L)$ with a vector $w\in\Lambda\cap L$.
Then $x-z-w\in bK+aK=(a+b)K$. Let $a,b$ decrease to the radii.
\end{proof}

\begin{proposition}[Previously established non-hollow bound]\label{base:nonhollow}
Every non-hollow lattice $d$-polytope satisfies $\mu(P,\Lambda)\le d/2$.
\end{proposition}
\begin{proof}
Translate an interior lattice point to zero and choose a minimal set
of vertices whose convex hull contains zero. Its hull is a simplex
$S$ of dimension $r\ge1$ with zero in its relative interior.
Every primitive length of its lattice vertices is at least one,
so Proposition~\ref{base:simplex} gives $\mu(S)\le r/2$.
If $r=d$, inclusion proves the result. Otherwise its rational span
$L$ has section radius at most $r/2$, and the quotient $\pi P$ is
a non-hollow lattice polytope of dimension $d-r$. Induction and
Lemma~\ref{base:section} give $r/2+(d-r)/2$.
\end{proof}

\begin{proposition}[Established weighted terminal radius]\label{base:weighted}
For positive real weights,
\[
 \mu(S_d(\omega),\Z^d)=
 B(c):=\frac{\sum_{j<k}c_jc_k}{\sum_jc_j},
 \qquad c_j=1/\omega_j.
\]
\end{proposition}
\begin{proof}
This is the formula of \cite[Theorem 1.4]{CSS}. For completeness,
the directions $-\mathbf1_d,e_1,\ldots,e_d$ generate $\Z^d$ and
have relation $(1,\ldots,1)$. The same coefficient minimization
as in Lemma~\ref{eq:deep}, with weights $c_j$, gives the exact
distance $\min_s\sum_jc_j\{t_j+s\}$. Its upper bound from
Lemma~\ref{base:flow} is $B(c)$. For the reverse inequality,
arrange its $d+1$ events in a cyclic order with a gap
$c_j/\sum_kc_k$ preceding the jump of weight $c_j$. These gaps
sum to one. Every rise is cancelled by the following jump, so
all post-jump values coincide. The jump identity gives their
common value $B(c)$, proving equality.
\end{proof}


\begin{lemma}[Exact flow distance in the vertex lattice]\label{base:exactflow}
Let $u_0,\ldots,u_r$ span an $r$-dimensional space and have positive
primitive relation $q$. Put $\Gamma=\sum_j\mathbb Zu_j$ and
$S=\operatorname{conv}(\lambda_0u_0,\ldots,\lambda_ru_r)$,
$\lambda_j>0$. For $x=\sum_jt_ju_j$,
\[
 \rho_{S,\Gamma}(x)=
 \min_s\sum_j\lambda_j^{-1}\{t_j+sq_j\}.
\]
\end{lemma}
\begin{proof}
Rounding the coefficients down supplies the upper bound. Conversely,
every lattice vector is $\sum_jn_ju_j$, $n_j\in\mathbb Z$, and every
nonnegative coefficient representation of $x-\sum_jn_ju_j$ in the
directions $u_j$ has coefficients $t_j-n_j+sq_j$. Each such
coefficient is at least its fractional part. Dividing by $\lambda_j$
and summing proves the lower bound.
\end{proof}

\begin{lemma}[Simplex tightness]\label{base:tight}
Every simplex $S$ is tight for every lattice: if a convex body
$K$ strictly contains $S$, then $\mu(K)<\mu(S)$.
This also applies to nonlattice super-bodies.
\end{lemma}
\begin{proof}
This is \cite[Lemma 2.7]{CSS}; we recall the argument.
At any point last covered by $\mu(S)S$, each facet of the simplex
is needed: otherwise all covering translates through that point
would lie in the $r$ half-spaces through the point coming from
the other facets. Their $r$ normals are independent, so these
half-spaces cannot cover a neighborhood, a contradiction.
Equivalently, every facet has a lattice translate whose relative
interior contains the point.

If $K$ strictly contains $S$, some facet of $S$ has its relative
interior in the interior of $K$. Its needed translate then puts
every old last-covered point in the interior of $\mu(S)K$.
All other points were already interior-covered at that scale.
Compactness of the lattice torus and continuity of the lattice
distance therefore give $\mu(K)<\mu(S)$.
\end{proof}


\section{Arithmetic event deficits and new interior-point families}
\label{cr6d:section}

The transfer used here is concrete: the finite sawtooth distribution
and residue-permutation argument in corpus \#030
(Lemma \texttt{p10:gn:sawtooth}, also equation
\texttt{p11:gc:sawtooth}) is applied to the \emph{values at the jump
events} of the covering flow in corpus \#200.  No transcendence or
formal gamma-calculus assertion from \#030 is needed.  We prove the
finite arithmetic ingredient directly, including the endpoints.
The result is a quantitative improvement of the unweighted covering
bound, rather than a new name for the already established strictness
of that bound.

\begin{lemma}[A finite residue deficit]\label{cr6d:residue}
Let $N,Q\ge1$ be integers and $g=\gcd(N,Q)$.  Suppose that real
numbers $X_0,\ldots,X_{N-1}$ satisfy
\[
 X_r-X_0\equiv \frac{Qr}{N}\pmod{\mathbb Z}.
\]
Then
\[
 \sum_{r=0}^{N-1}\bigl(X_r-\min_sX_s\bigr)\ge\frac{N-g}{2}.
\]
Consequently the same lower bound holds with any common lower
bound for the $X_r$ in place of their minimum.
\end{lemma}
\begin{proof}
Choose $r_0$ at which the minimum is attained.  The nonnegative
numbers $X_r-X_{r_0}$ have fractional parts
$\{Q(r-r_0)/N\}$.  Writing $N=gN'$ and $Q=gQ'$, with
$\gcd(N',Q')=1$, these fractional parts run through
$0,1/N',\ldots,(N'-1)/N'$, each exactly $g$ times.
Their sum is $g(N'-1)/2=(N-g)/2$.  A nonnegative real number
is at least its fractional part, proving the assertion.
\end{proof}

\begin{theorem}[Arithmetic fractional-part descent]\label{cr6d:flow}
For positive integers $q_0,\ldots,q_r$, put
\[
 Q=\sum_{j=0}^r q_j,\qquad
 G(q)=\sum_{j=0}^r\gcd(Q,q_j).
\]
For every real phase vector $t$,
\begin{equation}\label{cr6d:flowbound}
 \min_{s\in\mathbb R}\sum_{j=0}^r\{t_j+q_js\}
 \le \frac{r-1}{2}+\frac{G(q)}{2Q}.
\end{equation}
This includes coincident jump events.
\end{theorem}
\begin{proof}
Write $F(s)=\sum_j\{t_j+q_js\}$ and $m=\min_sF(s)$.
Let $E$ be its distinct jump events in one period, and let
$J_e$ be the number of coordinates jumping at $e$.  Values at
events are the right-continuous, post-jump values.
The square-balance identity of corpus \#200 gives
\[
 \sum_{e\in E}J_eF(e)
 =\frac{Q(r+1)}2-\frac12\sum_{e\in E}J_e^2,
 \qquad \sum_{e\in E}J_e=Q.
\]
Therefore
\begin{equation}\label{cr6d:balance}
 Q\left(\frac r2-m\right)
 =\sum_{e\in E}J_e(F(e)-m)
   +\frac12\sum_{e\in E}J_e(J_e-1).
\end{equation}

Fix coordinate $j$.  Its $q_j$ events may be parametrized by
$e_{j,b}=e_{j,0}+b/q_j$, for $0\le b<q_j$, with time taken
modulo one.  Since
$F(s)=\sum_i t_i+Qs-\sum_i\lfloor t_i+q_is\rfloor$,
\[
 F(e_{j,b})-F(e_{j,0})\equiv\frac{Qb}{q_j}\pmod{\mathbb Z}.
\]
Lemma~\ref{cr6d:residue}, and the fact that $m$ is a lower
bound for these event values, imply
\[
 \sum_{b=0}^{q_j-1}(F(e_{j,b})-m)
 \ge\frac{q_j-\gcd(Q,q_j)}2.
\]
Summing in $j$ counts a shared event exactly $J_e$ times, so
the first term on the right of \eqref{cr6d:balance} is at least
$(Q-G(q))/2$.  The second term is nonnegative.  Rearrangement
proves \eqref{cr6d:flowbound}.
\end{proof}

\begin{corollary}[An arithmetic covering certificate]\label{cr6d:simplex}
Let $S=\operatorname{conv}(v_0,\ldots,v_r)$ be a lattice
$r$-simplex with $0\in\operatorname{relint}S$, in its actual
ambient lattice.  If its primitive positive integral relation is
$\sum_jq_jv_j=0$, then
\begin{equation}\label{cr6d:radius}
 \mu(S)\le\frac{r-1}{2}+\frac{G(q)}{2Q}.
\end{equation}
Vertex generation of the ambient lattice is not required.

More generally, suppose a lattice $d$-polytope $P$ has $h\ge1$
interior lattice points, $0\in\operatorname{int}P$, and contains
such a simplex $S$ in a rational section through zero.  Then
\[
 \mu(P)\le\frac{d-1}{2}+\frac{G(q)}{2Q}.
\]
In particular, the numerical interior-point bound holds whenever
\begin{equation}\label{cr6d:criterion}
 (h+1)G(q)\le2Q.
\end{equation}
\end{corollary}
\begin{proof}
For any $x=\sum_jt_jv_j$, the vector
$z_s=\sum_j\lfloor t_j+q_js\rfloor v_j$ belongs to the lattice.
The relation gives
$x-z_s=\sum_j\{t_j+q_js\}v_j\in F(s)S$.
Theorem~\ref{cr6d:flow} proves \eqref{cr6d:radius}.
For the section statement, apply the section--projection
covering inequality.  The quotient is a non-hollow lattice
polytope of dimension $d-r$, so corpus \#200 bounds its covering
radius by $(d-r)/2$.  Adding proves the asserted bound for $P$.
\end{proof}

\begin{remark}[Rational weighted version]\label{cr6d:weighted}
If $c_j=\kappa n_j$ with $\kappa>0$ and positive integers $n_j$,
put $C'=\sum_jq_jn_j$, $C=\kappa C'$, and
\[
 B(q,c)=\frac12\sum_jc_j-\frac{\sum_jq_jc_j^2}{2C}.
 \]
The same proof, applied to $F/\kappa$ at coordinate $j$'s
events and then to the weighted square-balance identity, gives
\[
 \min_s\sum_jc_j\{t_j+q_js\}
 \le B(q,c)
 -\frac{\kappa}{2C}
   \sum_jc_j\bigl(q_j-\gcd(C',q_j)\bigr).
\]
Indeed $\sum_eJ_e(F(e)-m)=
\sum_jc_j\sum_b(F(e_{j,b})-m)$, while the cross terms in
$\sum_eJ_e^2-\sum_jq_jc_j^2$ are nonnegative.
This arithmetic refinement is asserted for the stated
commensurable weights; it is not an unproved extension to
arbitrary independent real weights.
\end{remark}

\subsection{Every non-hollow lattice simplex of prime normalized volume}

\begin{lemma}[An elementary interior-volume bound]\label{cr6d:volume}
If a lattice $d$-simplex has $h$ interior lattice points, then
its normalized volume $V$ satisfies
\[
 V\ge dh+1.
\]
\end{lemma}
\begin{proof}
Start with the simplex as a triangulation, and insert its $h$
interior lattice points one at a time by stellar subdivision.
If a new point lies in the relative interior of an $s$-face of
the current triangulation, $1\le s\le d$, that face is internal
because the point lies in the interior of the original simplex.
For $s<d$ its link is a triangulated sphere of dimension
$d-s-1$ and has at least $d-s+1$ facets.  To see the latter
elementary bound, fix one facet: each of its codimension-one
faces has a distinct neighboring facet.  For a zero-dimensional
sphere the bound says that it has at least two points.
For $s=d$ there is one incident top-dimensional simplex.

Each incident top-dimensional simplex is replaced by $s+1$
such simplices, so the number increases by at least
$s(d-s+1)\ge d$.  Thus the final triangulation has at least
$dh+1$ top-dimensional simplices.  All their vertices are
lattice points, so each has normalized volume at least one.
Adding their volumes proves the bound.  Dimension one is the
same argument with subintervals.
\end{proof}

\begin{theorem}[Coprime barycentric numerators]\label{cr6d:coprime}
Let $S$ be a lattice $d$-simplex with $h\ge1$ interior lattice
points, and translate one of them to zero.  Suppose that its
vertices generate the ambient lattice and that the primitive
positive relation $q$ satisfies
\[
 \gcd(q_j,Q)=1 \quad\hbox{for every }j,\qquad Q=\sum_jq_j.
\]
Then
\begin{equation}\label{cr6d:coprimebound}
 \mu(S)\le\frac{d-1}{2}+\frac{d+1}{2Q}
 \le\frac{d-1}{2}+\frac1{h+1}.
\end{equation}
If $d\ge2$ and $h\ge2$, the last inequality is strict.
\end{theorem}
\begin{proof}
The first bound follows from $G(q)=d+1$.
For a vertex-generated lattice simplex the normalized volume
is $Q$; the cofactor relation gives the normalized volumes of
the cones from zero as $q_0,\ldots,q_d$, and these cones
partition the simplex.  Lemma~\ref{cr6d:volume} gives
$Q\ge dh+1$.  Finally
\[
 2(dh+1)-(d+1)(h+1)=(d-1)(h-1)\ge0.
\]
This proves the second bound and the claimed strictness.
\end{proof}

\begin{corollary}[Prime normalized volume]\label{cr6d:prime}
Every non-hollow lattice simplex of prime normalized volume $p$
satisfies the numerical interior-point conjecture.  More precisely,
\[
 \mu(S)\le\frac{d-1}{2}+\frac{d+1}{2p}
 \le\frac{d-1}{2}+\frac1{h+1}.
\]
The comparison with the conjectured bound is strict when
$d\ge2$ and $h\ge2$.  No collinearity assumption on the interior
lattice points is made.
\end{corollary}
\begin{proof}
After translating an interior lattice point to zero, let
$I$ be the index of the vertex-generated lattice.  The cofactor
relation gives $p=IQ$, where $Q=\sum_jq_j\ge d+1>1$.
Primality forces $I=1$ and $Q=p$.  Each $q_j$ is strictly
between zero and $p$, so it is coprime to $p$.
Apply Theorem~\ref{cr6d:coprime}.
\end{proof}

\begin{example}[Three noncollinear interior points]\label{cr6d:example}
Consider
\[
 S=\operatorname{conv}\bigl(-(1,3,6),e_1,e_2,e_3\bigr).
\]
Its primitive relation is $(1,1,3,6)$, of sum $11$.
Its interior lattice points are exactly
\[
 (0,0,0),\qquad (0,0,-1),\qquad (0,-1,-2),
\]
which are noncollinear.  To check the count directly, the first
coordinate of an interior lattice point is zero.  Write the
other coordinates as $-u,-v$ with nonnegative integers $u,v$.
The barycentric inequalities become
$3+3v>8u$ and $6+6u>5v$.
For $u=0$ they give $v=0,1$; for $u=1$ they give $v=2$;
and $u\ge2$ is impossible.  The new certificate gives
\[
 \mu(S)\le1+\frac2{11}=\frac{13}{11}
 <\frac54=\frac{3-1}{2}+\frac1{3+1}.
\]
All four displayed vertices are primitive.  Thus this example
does not enter through the earlier certificate requiring
three nonprimitive vertices.
\end{example}

\subsection{An infinite family inside the old two-point gap}

\begin{theorem}[Two repeated relation weights]\label{cr6d:family}
Let $k,\ell,a$ be positive integers, $d=k+\ell-1$, and
\[
 S_{k,\ell,a}=
 \operatorname{conv}\bigl(-(1^{\,k-1},a^{\,\ell}),
                         e_1,\ldots,e_d\bigr)
 \subset\mathbb R^d .
\]
Here the superscripts indicate repeated coordinates, including
an empty block when $k=1$.  Set $g=\gcd(k,a)$ and
$h=\lceil a/k\rceil$.  Then
\begin{align}
 |\operatorname{int}S_{k,\ell,a}\cap\mathbb Z^d|&=h,
 \label{cr6d:count}\\
 \mu(S_{k,\ell,a})&=
 \frac{d-1}{2}+\frac{k+\ell g}{2(k+a\ell)}
 \le\frac{d-1}{2}+\frac1{h+1}.
 \label{cr6d:familybound}
\end{align}
For $h\ge2$ equality with the conjectured bound holds exactly
when $\ell=1$ and $k\mid a$; in particular, it is strict for
$\ell\ge2$.  For $h=1$, equality with $d/2$ holds exactly when
$a\mid k$.
\end{theorem}
\begin{proof}
The relation has $k$ entries equal to one and $\ell$ entries
equal to $a$, so $Q=k+a\ell$ and
$G(q)=k+\ell\gcd(k,a)$.  This proves the asserted upper bound.
We next construct phases attaining it, so the radius formula
is exact.

Write $\theta_b=\{kb/a\}$ and
\[
 N_b=\left\lfloor\frac{(b+1)k}{a}\right\rfloor
     -\left\lfloor\frac{bk}{a}\right\rfloor
 \qquad (0\le b<a).
\]
Schedule the events of the $\ell$ frequency-$a$ coordinates at
\[
 e_{b,j}=\frac ba+\frac jQ,
 \qquad 0\le b<a,\quad 0\le j<\ell.
\]
For each fixed $j$ these are equally spaced by $1/a$, so they
are the events of a genuine frequency-$a$ coordinate.
In the gap after the block with index $b$, place $N_b$
frequency-one events at
\[
 u_{b,p}=\frac ba+\frac{\ell-1+p-\theta_b}{Q},
 \qquad 1\le p\le N_b .
\]
Their total number is $\sum_bN_b=k$.  The first is strictly
after $e_{b,\ell-1}$ because $1-\theta_b>0$, and the last is
strictly before $e_{b+1,0}$ because the remaining gap is
$(1+\theta_{b+1})/Q>0$.  If $N_b=0$ there are no intervening
unit events.  Thus this is a simple, noncoincident event
schedule of the required coordinates.

Let $M$ denote its post-jump value at $e_{0,0}=0$.
The slope of the flow is $Q$.  Consecutive events within each
frequency-$a$ block are separated by $1/Q$ and have jumps
one, so their post-jump values agree.  An induction on $b$
shows that their common value is $M+\theta_b$.
Every intervening unit event has post-jump value $M$:
the waiting time to its $p$th event, measured from the last
event of the block, is $(p-\theta_b)/Q$, and exactly $p$
unit jumps have then occurred.  The next block has value
$M+\theta_{b+1}$ because
$\theta_{b+1}=\theta_b+k/a-N_b$.
Consequently the minimum of this flow is $M$.

All jumps are one.  The square-balance identity gives
\[
 \frac{Qd}{2}
 =\sum_{\text{all events }e}F(e)
 =QM+\ell\sum_{b=0}^{a-1}\theta_b
 =QM+\frac{\ell(a-g)}2.
\]
Hence
\[
 M=\frac d2-\frac{\ell(a-g)}{2Q}
   =\frac{d-1}{2}+\frac{k+\ell g}{2Q}.
\]
The vertices of $S_{k,\ell,a}$ generate $\mathbb Z^d$.
For completeness, for a point $x=\sum_jt_jv_j$ and any
lattice vector $z=\sum_jn_jv_j$, every nonnegative
coefficient vector representing $x-z$ has the form
$t-n+sq$, since the kernel of the vertex map is the line
spanned by $q$.  Its coefficient sum is at least
$\sum_j\{t_j+sq_j\}$.  The preceding rounding argument
gives the converse inequality.  Thus the distance of $x$
to the lattice is exactly the minimum of its flow.
The constructed phases therefore attain $M$, proving the
radius equality.

An interior lattice point has zero in its first $k-1$
coordinates, since each lies strictly between $-1$ and $1$.
Write its other coordinates as $-n_1,\ldots,-n_\ell$ with
nonnegative integers $n_j$.  Its coefficient at the negative
vertex is $(1+\sum_jn_j)/Q$.  Put
$M=\max_jn_j$ and $D=\ell M-\sum_jn_j$.
Positivity of the coefficient at a coordinate realizing $M$
requires
\[
 a(1-D)>kM.
\]
Hence $D=0$: all $n_j$ equal $M$.  The remaining condition is
$kM<a$.  The interior points are therefore exactly
$(0^{k-1},-M^\ell)$ for $0\le M<h$, proving
\eqref{cr6d:count}.

For $h=1$ the desired bound is the already proved $d/2$
bound, and also follows from $g\le a$ in
\eqref{cr6d:familybound}.  Suppose $h\ge2$.
Write $k=gK$, $a=gA$, and
$A=K(h-1)+R$ with $1\le R\le K$.  Then
\[
 \begin{split}
 2(K+\ell A)-(h+1)(K+\ell)
 &=(h-1)\bigl((2\ell-1)K-\ell\bigr)
       +2\ell(R-1)\\
 &=(h-1)\bigl((\ell-1)K+\ell(K-1)\bigr)
       +2\ell(R-1)\ge0.
 \end{split}
 \]
This is precisely the comparison with the conjectured bound in
\eqref{cr6d:familybound}.  When $h\ge2$, equality in this
nonnegative expression forces $\ell=1$, $K=1$, and $R=1$.
These conditions are equivalent to $\ell=1$ and $k\mid a$.
For $h=1$, the exact radius formula equals $d/2$ if and only
if $g=a$, equivalently $a\mid k$.
\end{proof}

\begin{remark}[Why this reaches genuinely unhandled two-point cases]
For $k<a\le2k$, the family has exactly two interior lattice
points.  Along their line, with primitive direction
$u=(0^{k-1},-1^\ell)$, its section is
\[
 S_{k,\ell,a}\cap\mathbb Ru
 =[-\ell^{-1}u,(a/k)u].
\]
Thus the notation of the earlier line method is
$\beta=1/\ell$, $\delta=(a-k)/k$.  With
$a=k+1$, $\ell=2$, and every $k\ge2$,
\[
 2\max(\beta,\delta)=1
 <1+\min(\beta,\delta)=1+\frac1k .
\]
These are exactly within the previously unhandled endpoint
regime.  Theorem~\ref{cr6d:family} now proves the conjectured
inequality throughout this infinite family, in dimensions
$d=k+1$.  It gives the explicit stronger estimate
\[
 \mu(S_{k,2,k+1})
 = \frac{k}{2}+\frac{k+2}{2(3k+2)}.
\]
For example, $k=4$, $\ell=2$, $a=5$ gives a five-dimensional
simplex of normalized volume $14$ with exactly two interior
lattice points and
\[
 \mu(S_{4,2,5})=\frac{31}{14}<\frac73
 =\frac{5-1}{2}+\frac1{2+1}.
\]
\end{remark}

\begin{remark}[The bridge that remains open]\label{cr6d:limit}
The arithmetic certificate \eqref{cr6d:criterion} need not
hold for every simplex with $h$ interior points.  For example,
the lattice triangle with relation $(1,3,4)$ has $Q=8$,
$h=2$, and $G=1+1+4=6$.  The resulting bound is $7/8$,
which is larger than the conjectured $5/6$.
The planar conjecture is known by other arguments, so this
is a failure of the present certificate, not a counterexample.
Neither the unrestricted numerical conjecture nor the whole
remaining two-point endpoint regime is resolved by this section.
\end{remark}


% New weighted result for the sixth creative-transfer round.
% This is an include fragment. Standard notation: \R, \Z, \conv, \mu.
% No historical-priority claim is made.

\section{Three distinct weights: an exact second covering minimum}
\label{c6w:sec}

Let
\[
 S_n(\omega_0,\ldots,\omega_n)
 =\conv\{-\omega_0\mathbf1_n,\omega_1e_1,\ldots,\omega_ne_n\}.
\]
The previous equal-pair and one-small-weight formulas do not determine
the case of three distinct weights.  The result below closes an infinite
family of that case, with arbitrary real overall scale and final weight.

\begin{theorem}[Three real weights in every dimension]\label{c6w:main}
Let $n\ge2$ and let $0<a\le b$, $c\ge2b$ be arbitrary real
numbers.  Put $A=1/a$, $B=1/b$, $C=1/c$, $T=A+B$, and
\[
 \varepsilon_k=\min\left\{
 A\left\{\frac{kB}{T}\right\},\;
 B\left(1-\left\{\frac{kB}{T}\right\}\right)\right\}.
\]
Then
\begin{equation}\label{c6w:answer}
 \boxed{\displaystyle
 \mu_2\bigl(S_n(a,b,c,\ldots,c),\Z^n\bigr)
 =\max_{1\le k\le n-1}
 \left\{\frac{AB}{T}
 +\frac{C(T-\varepsilon_k)}{T+kC}\right\}.}
\end{equation}
Every term in the maximum is attained by an explicit rational
rank-two quotient.  No integrality or rationality assumption is
made on $a,b,c$.
\end{theorem}

\begin{corollary}[A closed expression for integer weight ratios]
\label{c6w:ratio}
Let $m\ge2$ be an integer, $a=b/m$, and $c\ge2b$.
For every $n\ge2$ put $h=\min\{n-1,m+1\}$.  Then
\begin{equation}\label{c6w:ratioanswer}
 \boxed{\displaystyle
 \mu_2\bigl(S_n(a,b,c,\ldots,c)\bigr)
 =\max\left\{
 \frac{a+b+c}{ab+ac+bc},\;
 \frac{c+b+ha}{c(a+b)+hab}
 \right\}.}
\end{equation}
The first expression wins when $2b\le c\le mb$; the second
wins when $c\ge mb$.  If $n\ge m+2$, the second expression is
\[
 \frac1{a+b}+\frac1{b+c}.
\]
The same answer holds if the weights with indices
$2,\ldots,h+1$ are kept equal to $c$ and all later weights are
enlarged arbitrarily.  In the coordinate branch $2b\le c\le mb$,
only the weight with index $2$ needs to remain equal to $c$.
\end{corollary}

For example,
\[
 \mu_2\bigl(S_n(1,2,6,\ldots,6)\bigr)=\frac{11}{24}
 \qquad(n\ge4).
\]
The coordinate projection gives only $9/20$; the exact excess is
$1/120$.  Thus the new formula is not obtained by substituting into
the original coordinate-projection conjecture.

The proof combines the jump identity of VMRC \#200 with a finite
placement of unequal events and a circular pigeonhole argument.
VMRC \#029, Lemma \texttt{p03:bn:returns} (``Simultaneous returns
of finitely many phases''), supplies the elementary torus
pigeonhole mechanism; the one-dimensional finite form needed
here is proved explicitly below, with a bound tied to $n$.
The methodological comparison with
VMRC \#215 is its reduction from continuous angles to a finite
weighted partition.  Here the continuous phase problem is reduced
to two item placements on a cyclic grid.  The circular pigeonhole
step bounds the needed multiplicity by the ambient dimension; for
integer weight ratios a residue period then removes the finite
maximum as well.  No theorem
about matrix spectra is used as a covering theorem.

\subsection{Two unequal events on a cyclic grid}

For positive real numbers $A\ge B$ and $C>0$, put $T=A+B$ and,
for each integer $k\ge1$, define
\begin{align}
 x_k&=\frac{kB}{T},\qquad u_k=\{x_k\},\nonumber\\
 \varepsilon_k&=\min\{A u_k,\, B(1-u_k)\},\label{c6w:epsilon}\\
 R_k(A,B,C)&=\frac{AB}{T}
       +\frac{C(T-\varepsilon_k)}{T+kC}.\label{c6w:Rk}
\end{align}
When $u_k=0$, formula \eqref{c6w:epsilon} gives
$\varepsilon_k=0$.

\begin{lemma}[Exact two-event placement]\label{c6w:flow}
For all $k\ge1$,
\begin{equation}\label{c6w:floweq}
 \max_{t_0,t_1,t_2\in\R}\min_{s\in\R}
 \bigl(A\{t_0+s\}+B\{t_1+s\}+C\{t_2+ks\}\bigr)
 =R_k(A,B,C).
\end{equation}
Equivalently, the triangle
\[
 P_k(a,b,c)=\conv\{(-ka,-a),(0,b),(c,0)\},
 \qquad (A,B,C)=(a^{-1},b^{-1},c^{-1}),
\]
has covering radius $R_k(A,B,C)$ with respect to $\Z^2$.
\end{lemma}

\begin{proof}
We spell out the optimization of the previously established
ordered-partition identity.  The identity follows by recording
individual jumps, including a fixed order for simultaneous jumps:
if $D=T+kC$ is the slope, $M$ is the minimum, and $H_e$ the
formal post-jump height at a jump of size $w_e$, then
\begin{equation}\label{c6w:deficit}
 M=\mathcal B-\frac1D\sum_e w_e(H_e-M),
 \qquad
 \mathcal B=\frac{AB+(k+1)CT/2}{T+kC}.
\end{equation}
All post-jump excesses are nonnegative.  At coincident events the
formal intermediate state is included, so that its contribution
accounts for the cross term in the square of the combined jump.

Rotate time so that the $C$-jumps occur at $h/k$ for
$h=0,\ldots,k-1$.  Cyclically label the interval containing the
$A$-jump as interval zero and the interval containing the $B$-jump
as interval $j$, where $0\le j\le k-1$.  If both lie in the same
interval, put $A$ before $B$: an adjacent interchange shows that
decreasing order of jump sizes minimizes the deficit.  Indeed,
for $A\ge B$ and any level $X$,
\[
 A(X-A)_++B(X-A-B)_+
 \le B(X-B)_++A(X-A-B)_+.
\]

Let $D_h=hT/k-\sum_{r<h}W_r$, where $W_r$ is the total singleton
weight assigned to interval $r$, and set
$Y_h=D_h-\min_{0\le r<k}D_r$.
The least possible deficit for this assignment is
\begin{equation}\label{c6w:cost}
 \mathcal T_j=C\sum_{h=0}^{k-1}Y_h
 +\sum_{\text{singleton jumps }w\text{ in a bin}}
      w(Y_h-\text{prefix sum through }w)_+.
\end{equation}
For completeness, these are lower bounds on every post-jump
excess because elapsed time only raises the level.  They can all
be attained: beginning a bin at excess $Y_h$, prescribe successive
excesses $Z_r=(Y_h-\text{prefix sum through }r)_+$ and wait
$(w_r+Z_r-Z_{r-1})/D\ge0$ before each singleton.  Before the next
distinguished jump wait $(C+Y_{h+1}-Z_{\rm last})/D>0$.
The total waiting time is
$(W_h+C+Y_{h+1}-Y_h)/D=1/k$.
At least one distinguished level is zero, and the prescription
therefore realizes an actual minimum.  These events are exactly
the jumps of the three fractional-part coordinates.  Their mean,
or \eqref{c6w:deficit}, determines the additive baseline.  Thus
minimizing \eqref{c6w:cost} gives the exact maximum in
\eqref{c6w:floweq}, rather than merely an upper bound.

For $k=1$, the deficit vanishes and the claimed answer follows
immediately.  Suppose $k\ge2$.  Directly summing the two rising
arithmetic progressions in the $Y_h$ gives
\begin{equation}\label{c6w:area}
 \sum_hY_h=\frac{T(k-1)}2-\min\{Aj,B(k-j)\}.
\end{equation}
This also holds for $j=0$, when both singleton events occupy the
same interval.  More explicitly, for $1\le j\le k-1$, the two
possible minima of $D_h$ are $T/k-A$ and $(j+1)T/k-T$;
their difference is $B-jT/k$.

The right side of \eqref{c6w:area} is minimized by one of
$j=\lfloor x_k\rfloor,\lceil x_k\rceil$.  Both choices lie in
$\{0,\ldots,k-1\}$, except for a redundant endpoint when $k=1$.
For these nearest choices all singleton terms in
\eqref{c6w:cost} vanish.  If $j\le x_k$, the only potentially
nonzero term is $A(B-(j+1)T/k)_+$; if $j\ge x_k$, it is
$B((j-1)T/k-B)_+$.  When $j=0$ the first expression is the
same-bin term.  Each is zero at the appropriate nearest integer.
Consequently
\[
 \min_j\mathcal T_j
 =C\left(\frac{T(k-1)}2-M_k\right),
 \quad
 M_k=\max_{0\le j\le k-1}\min\{Aj,B(k-j)\}.
\]
The two nearest integers show
\[
 M_k=\frac{kAB}{T}-\varepsilon_k.
\]
Substitution into \eqref{c6w:deficit} yields
\[
 R_k=\frac{AB+C(T+M_k)}{T+kC}
 =\frac{AB}{T}+\frac{C(T-\varepsilon_k)}{T+kC}.
\]
Finally the integer directions $(-k,-1),(0,1),(1,0)$ generate
$\Z^2$ and have primitive positive relation $(1,1,k)$.
The exact gauge identity therefore identifies this flow optimum
with the asserted covering radius.
\end{proof}

\subsection{Eliminating all other planar relations}

We use the previously proved one-small-column planar result in
the following precise form.  For lattice directions surrounding
zero, one column scaled by $a$ and every other column scaled by
at least $c\ge a$ give radius at most
\[
 \frac{C(2A+C)}{A+2C},\qquad A=1/a,\ C=1/c.
\]
Its proof is the planar unimodular-cone refinement and the
three first-event inequality.  This already established result is
only used below for triangles having at most one small column.
Its full proof is reproduced in Appendix~\ref{app:onesmall}.

\begin{theorem}[Sharp planar envelope]\label{c6w:envelope}
Let $0<a\le b\le c$, with $c\ge2b$, and write
$A=1/a$, $B=1/b$, $C=1/c$.  Suppose finitely many vectors of a
planar lattice surround zero.  Scale at most one vector by $a$,
at most one further vector by $b$, and all the others by real
factors at least $c$.  For their convex hull $P$,
\begin{equation}\label{c6w:envelopeeq}
 \mu(P)\le \max_{k\ge1}R_k(A,B,C).
\end{equation}
The maximum exists and the bound is attained within this class
by one of the triangles $P_k(a,b,c)$.
\end{theorem}

\begin{proof}
Shrink all larger multipliers to $c$.  Ignore zero columns when
choosing a positive dependence of minimal support; removing
them does not affect origin-interiority of the convex hull of
the remaining columns.  In the plane the chosen support is
either an origin-containing triangle or an opposite pair.

First suppose the triangle uses both small columns.  Let
$(p,q,\ell)$ be the primitive positive integral relation between
the unscaled directions, with $p$ attached to weight $A$ and
$q$ to weight $B$.  The usual rounding argument bounds its
covering radius by the maximum over phases of
\[
 \min_s\bigl(A\{t_0+ps\}+B\{t_1+qs\}
                         +C\{t_2+\ell s\}\bigr).
\]
This upper bound does not require the three directions to
generate the ambient lattice.

For positive integers $p,q$ and $g=\gcd(p,q)$, the exact
two-coordinate bound is
\begin{equation}\label{c6w:twofreq}
 \min_s\bigl(A\{t_0+ps\}+B\{t_1+qs\}\bigr)
 \le\frac{gAB}{pA+qB}.
\end{equation}
Indeed the two integer directions $q/g,-p/g$ generate $\Z$;
the corresponding interval has length
$q/(gA)+p/(gB)$, and its covering radius is the reciprocal
of that length.

If $p\ne q$, then $p/g,q/g$ are unequal positive integers.
As $A\ge B$,
\[
 \mu\le\frac{AB}{A+2B}+C\le R_2(A,B,C),
\]
because
\[
 R_2-\left(\frac{AB}{A+2B}+C\right)
 =\frac{(B-2C)(AB+AC+2BC)}
 {(A+2B)(A+B+2C)}\ge0.
\]
If $p=q=g\ge2$, primitivity implies $\gcd(g,\ell)=1$.
Choose a time giving the first two terms at most $AB/(A+B)$.
The shifts by $j/g$ preserve these terms and send the last
fractional part through a complete $g$-point grid.  One value
is at most $1/g$.  Therefore
\[
 \mu\le\frac{AB}{A+B}+\frac C2\le R_1(A,B,C),
\]
where the last difference is
\[
 \frac{C(A^2+B^2-C(A+B))}{2(A+B)(A+B+C)}\ge0.
\]
The only remaining case is $p=q=1$, precisely the exact family
of Lemma~\ref{c6w:flow}.

Write
\[
 E=R_1(A,B,C)=\frac{AB+AC+BC}{A+B+C}.
\]
A triangle with no small column has radius at most $C\le E$.
A triangle with only one small column has radius at most
$E(A,C,C)$ or $E(B,C,C)$ by the one-small-column result.
The function $(XY+XZ+YZ)/(X+Y+Z)$ increases in each positive
variable, so both bounds are at most $E(A,B,C)$.

For an opposite pair, use section/projection subadditivity.
If the pair contains both small columns, or neither, the sum of
the one-dimensional bounds is at most
\[
 \frac{AB}{A+B}+\frac C2\le E.
\]
If the pair contains exactly one small column, that sum is at
most
\[
 \frac{AC}{A+C}+\frac{BC}{B+C}\le E,
\]
since the difference in this case is
\[
 \frac{AB(AB-C^2)}{(A+C)(B+C)(A+B+C)}\ge0.
\]
These section lengths follow from primitivity of the line
coordinate: nonzero projected lattice columns are integers,
so the selected opposite pair has the sum of its two radial
multipliers as a lower bound, and the quotient interval has
the sum of the two smallest remaining multipliers as a lower
bound.  Origin-interiority guarantees both signs in the quotient.

We have covered every minimal positive dependence.  Finally,
$R_k\to AB/(A+B)<R_1$ as $k\to\infty$, so the maximum is attained.
The triangle $P_k$ attaining it belongs to the stipulated class,
and Lemma~\ref{c6w:flow} proves sharpness.
\end{proof}

\subsection{The integer ratio removes the infinite maximum}

\begin{lemma}[Residue-period reduction]\label{c6w:period}
If $A=mB$ for an integer $m\ge1$, then for every $C>0$,
\begin{equation}\label{c6w:periodmax}
 \max_{k\ge1}R_k(A,B,C)
 =\max\left\{
 \frac{AB+AC+BC}{A+B+C},\;
 \frac{AB}{A+B}+\frac{BC}{B+C}
 \right\}.
\end{equation}
The two values are attained at $k=1$ and $k=m+1$, respectively.
The second wins precisely when $C\le B/m$.
\end{lemma}

\begin{proof}
Now $x_k=k/(m+1)$, so $\varepsilon_k$ is periodic with period
$m+1$.  Write $k=t(m+1)+r$, with $1\le r\le m+1$.  The numerator
$C(T-\varepsilon_k)$ in \eqref{c6w:Rk} is the same as for $r$
and strictly positive, whereas its denominator is no smaller.
Thus $R_k\le R_r$.

For $1\le r\le m+1$, the optimization in the flow lemma gives
$M_r=B(r-1)$ and hence
\[
 R_r=\frac{AB+C(A+rB)}{T+rC}.
\]
Consecutive differences are
\[
 R_{r+1}-R_r
 =\frac{C(B^2-AC)}{(T+rC)(T+(r+1)C)}.
\]
This finite sequence is monotone.  Its endpoints are $R_1=E$
and
\[
 R_{m+1}=\frac{AB}{A+B}+\frac{BC}{B+C}.
\]
Since $A=mB$, the latter is larger exactly when $C<B/m$;
all members coincide at equality.
\end{proof}

\subsection{The ambient dimension bounds the multiplicity}

The sharp planar envelope allows arbitrary lattice triangles.
Quotients of $S_n$ carry additional information: their unscaled
columns sum to zero.  Keeping this information is the step that
turns the envelope into an exact answer in \emph{every} dimension.

\begin{lemma}[A circular discrepancy bound]\label{c6w:circle}
For every integer $N\ge1$, at least one $j\in\{1,\ldots,N\}$
satisfies
\[
 \varepsilon_j\le\frac{A}{N+1}<\frac{T}{N+1}.
\]
\end{lemma}
\begin{proof}
Place the $N+1$ points
$0,\{B/T\},\ldots,\{NB/T\}$ on the unit circle.  A repeated
point immediately gives some $\varepsilon_j=0$.  Otherwise
some cyclic gap has length at most $1/(N+1)$.  Its endpoint
indices differ by an integer $j$ with $1\le j\le N$; thus
either $\{jB/T\}$ or $1-\{jB/T\}$ is at most $1/(N+1)$.
Since $B\le A$, the asserted bound follows from the definition
of $\varepsilon_j$.  The strict inequality uses $B>0$.
\end{proof}

\begin{proof}[Proof of Theorem~\ref{c6w:main}]
For $n=2$, this is the known full-dimensional weighted-simplex
formula $\mu_2=R_1$.
Suppose $n\ge3$ and put $N=n-1$.  Take an arbitrary rational
rank-two quotient $L$ and its lattice $\Lambda=L\Z^n$.  Write
\[
 u_0=L(-\mathbf1_n),\qquad u_j=Le_j\quad(1\le j\le n).
\]
These columns span the plane and satisfy
\begin{equation}\label{c6w:zerosum}
 u_0+u_1+\sum_{j=2}^n u_j=0.
\end{equation}
Their convex hull therefore has zero in its interior.

Use the minimal-positive-dependence classification from the
proof of Theorem~\ref{c6w:envelope}, ignoring zero columns only
for that choice.  Retain all original $N$ large columns,
including any zero ones, in \eqref{c6w:zerosum} and in the
average below.  Every opposite-pair case,
every triangle with at most one small column, and every triangle
whose primitive relation has $p\ne q$ or $p=q\ge2$, is already
bounded by $R_1$ or $R_2$.  Since $N\ge2$, both are allowed
terms in the desired finite maximum.

The remaining triangle uses $u_0,u_1,u_i$ with
\[
 u_0+u_1+\ell u_i=0
\]
for an integer $\ell\ge1$.  If $\ell\le N$, its radius is at
most $R_\ell(A,B,C)$ by the rounding argument and
Lemma~\ref{c6w:flow}.  Suppose instead that $\ell>N$.
Equation \eqref{c6w:zerosum} yields
\[
 \frac1N\sum_{j=2}^n u_j=\frac\ell N u_i.
\]
The quotient body therefore contains the larger triangle
\[
 \conv\left\{a u_0,\,b u_1,\,\frac{c\ell}{N}u_i\right\}.
\]
The same relation remains $(1,1,\ell)$, and the reciprocal
weight at its third vertex is now $CN/\ell$.  Consequently
its radius is at most
\begin{equation}\label{c6w:enlarged}
 R_\ell(A,B,CN/\ell)
 =\frac{AB}{T}
  +\frac{CN(T-\varepsilon_\ell)}{\ell(T+NC)}.
\end{equation}
Choose $j\le N$ from Lemma~\ref{c6w:circle}.  Since
$\ell\ge N+1$ and $\varepsilon_\ell\ge0$,
\[
 \frac{N(T-\varepsilon_\ell)}\ell
 \le\frac{NT}{N+1}
 <T-\varepsilon_j.
\]
Using $T+jC\le T+NC$ in \eqref{c6w:enlarged} now gives
\[
 R_\ell(A,B,CN/\ell)
 <\frac{AB}{T}+\frac{C(T-\varepsilon_j)}{T+NC}
 \le R_j(A,B,C).
\]
Thus no multiplicity larger than $N$ is needed.  We have proved
the required upper bound in every rational rank-two quotient;
the quotient characterization of covering minima proves the
upper bound in \eqref{c6w:answer}.

For each $1\le k\le N$, the surjective integer map
\[
 L_k(x_1,\ldots,x_n)
 =\left(\sum_{j=2}^{k+1}x_j,\,x_1\right)
\]
maps the simplex exactly onto $P_k(a,b,c)$; all further columns
map to zero.  Its exact radius is $R_k$ by
Lemma~\ref{c6w:flow}.  Taking the largest of these quotient
radii proves the matching lower bound.
\end{proof}

\begin{proof}[Proof of Corollary~\ref{c6w:ratio}]
If $A=mB$, the residue-period argument gives $R_k\le R_r$
whenever $k=t(m+1)+r$, $1\le r\le m+1$.
Thus the finite maximum over $k\le n-1$ reduces to
$1\le r\le h=\min\{n-1,m+1\}$.  On this initial interval,
\[
 R_r=\frac{AB+C(A+rB)}{A+B+rC},
\]
and the consecutive differences in the proof of
Lemma~\ref{c6w:period} show monotonicity with sign $B^2-AC$.
The maximum is therefore $\max\{R_1,R_h\}$.  Translating from
reciprocal weights gives \eqref{c6w:ratioanswer} and the branch
threshold $c=mb$.  If $h=m+1$, the endpoint is
$AB/(A+B)+BC/(B+C)$.

For the extension with enlarged later weights, the universal
planar envelope still supplies the same upper bound when
$n-1\ge m+1$.  When $n-1<m+1$, no later weights are present
in the stated extension.  The quotient $L_h$ keeps the same
image and supplies the matching lower bound.  In the coordinate
branch the global envelope is already $R_1$, and the coordinate
quotient alone supplies the lower bound, irrespective of the
later enlargements.
\end{proof}

\paragraph{Scope.}
The argument settles all three-weight families
$S_n(a,b,c,\ldots,c)$ with $c\ge2b$ and $0<a\le b$, in every
dimension.  It also settles the displayed cases with arbitrarily
enlarged later weights.  It does not settle arbitrary independent
weights, the range $b<c<2b$, or higher covering indices.  No
unrestricted weighted-conjecture resolution or historical-priority
claim is made.


% Fragment for the sixth research pass.  All statements below are proved here.
% The containing document needs amsmath, amssymb, amsthm and theorem environments.

\section{Covering products under arbitrary monotone gauge aggregation}
\label{c6f:sec:aggregation}

Write
\[
 \mathcal F_d(K,\Lambda)
   =\frac{\operatorname{vol}_d(K)}{\det\Lambda}
      \prod_{i=1}^d\mu_i(K,\Lambda).
\]
The classical covering-product conjecture asks for
$\mathcal F_d(K,\Lambda)\ge(d+1)/2^d$ for every convex body and
full-rank lattice.  We prove a closure principle for a class substantially
larger than Cartesian products and free sums.  It includes arbitrary
convex, coordinatewise monotone ways of combining the gauges of several
bodies.  The general conjecture is not asserted.

\subsection{The coefficient body and its weighted volume}

Let $m\ge1$, let $d_1,\ldots,d_m$ be positive integers, and put
$D=\sum_jd_j$.  A \emph{coefficient body} is a compact convex set
$C\subset\mathbb R_{\ge0}^m$ which contains $[0,\epsilon]^m$ for
some $\epsilon>0$ and is coordinatewise downward closed:
\[
 0\le u_j\le t_j\ (1\le j\le m),\quad t\in C
 \quad\Longrightarrow\quad u\in C.
\]
Its gauge on the positive orthant is
\[
 \gamma_C(t)=\inf\{\rho>0:t\in\rho C\}.
\]
It is continuous, positively homogeneous, convex, and coordinatewise
nondecreasing.  Define the weighted coefficient volume
\begin{equation}\label{c6f:eq:weighted-volume}
 W_{\boldsymbol d}(C)
   =\left(\prod_{j=1}^m d_j\right)
      \int_C\prod_{j=1}^m t_j^{d_j-1}\,dt_1\cdots dt_m.
\end{equation}

For a convex body $K_j\subset E_j$ containing the origin in its
interior, write $g_j$ for its Minkowski functional.  In the direct-sum
space $E_1\oplus\cdots\oplus E_m$, define
\begin{equation}\label{c6f:eq:aggregation}
 K_C=\{(x_1,\ldots,x_m):(g_1(x_1),\ldots,g_m(x_m))\in C\}.
\end{equation}
This is a convex body: convexity of each $g_j$, followed by downward
closure and convexity of $C$, proves convexity of the displayed set.
Compactness and a neighborhood of the origin follow from the
corresponding properties of the factors and of $C$.

Radial integration gives
\begin{equation}\label{c6f:eq:body-volume}
 \operatorname{vol}_D(K_C)
 =W_{\boldsymbol d}(C)\prod_{j=1}^m\operatorname{vol}_{d_j}(K_j).
\end{equation}
Indeed, the distribution function of $g_j$ under Lebesgue measure is
$\operatorname{vol}_{d_j}\{g_j\le t\}
=t^{d_j}\operatorname{vol}_{d_j}(K_j)$, so its pushforward measure
has density $d_jt^{d_j-1}\operatorname{vol}_{d_j}(K_j)$.

The choices
\[
 C=\{t\ge0:\textstyle\sum_jt_j\le1\},\qquad
 C=[0,1]^m,
\]
give the free sum and Cartesian product, respectively.  Choosing
$C=\{t\ge0:\sum_jt_j^r\le1\}$ gives the $\ell_r$ gauge sum.
The coefficient body need not have one of these special forms.

\subsection{A packing inequality for gauge convolutions}

\begin{lemma}[Permutation-simplex packing]\label{c6f:lem:packing}
For every $j$, let $a_{j,1},\ldots,a_{j,d_j}$ be positive real
numbers, and set $a_{j,0}=0$.  For $1\le k\le D$ define
\begin{equation}\label{c6f:eq:convolution}
 c_k=\max_{\substack{0\le i_j\le d_j\ (1\le j\le m)\\
                     i_1+\cdots+i_m=k}}
         \gamma_C(a_{1,i_1},\ldots,a_{m,i_m}).
\end{equation}
Then
\begin{equation}\label{c6f:eq:packing-bound}
 W_{\boldsymbol d}(C)\prod_{k=1}^D c_k
   \ge\prod_{j=1}^m\prod_{i=1}^{d_j}a_{j,i}.
\end{equation}
No monotonicity of the input sequences is needed.
\end{lemma}

\begin{proof}
Use $D$ auxiliary nonnegative coordinates, partitioned into blocks of
sizes $d_j$.  In block $j$, let
$x_{j,(1)}\ge\cdots\ge x_{j,(d_j)}\ge0$ denote the coordinates
in decreasing order, with $x_{j,(d_j+1)}=0$.  Set
\begin{align*}
 f_j(x_j)
  &=\sum_{i=1}^{d_j}a_{j,i}
           (x_{j,(i)}-x_{j,(i+1)})\\
  &=\sum_{i=1}^{d_j}(a_{j,i}-a_{j,i-1})x_{j,(i)}.
\end{align*}
Although $f_j$ need not be convex on the whole positive orthant, it
is linear on each region with a fixed coordinate order.  It is
continuous, positively homogeneous and strictly positive away from
zero, since the first expression uses positive $a_{j,i}$ and
nonnegative consecutive differences.

Let $Q_j=\{x_j\ge0:f_j(x_j)\le1\}$.  In the ordered region
$x_{j,1}\ge\cdots\ge x_{j,d_j}\ge0$, make the triangular change
of variables $y_i=x_{j,i}-x_{j,i+1}$.  It has determinant one and
turns $Q_j$ in this region into
$\{y\ge0:\sum_i a_{j,i}y_i\le1\}$.  Hence that region has
volume $1/(d_j!\prod_i a_{j,i})$.  The $d_j!$ coordinate-order
regions have disjoint interiors, so
\begin{equation}\label{c6f:eq:aux-volume}
 \operatorname{vol}_{d_j}(Q_j)=\frac1{\prod_i a_{j,i}}.
\end{equation}

Now form the auxiliary set
\[
 A=\{x\in\mathbb R_{\ge0}^{D}:
                (f_1(x_1),\ldots,f_m(x_m))\in C\}.
\]
The radial integration argument proving
\eqref{c6f:eq:body-volume} only needs homogeneity and the volume of
the unit sublevel set; it does not need global convexity.  Therefore
\begin{equation}\label{c6f:eq:A-volume}
 \operatorname{vol}_D(A)
   =\frac{W_{\boldsymbol d}(C)}
          {\prod_{j,i}a_{j,i}}.
\end{equation}

Fix a permutation $\sigma$ of all $D$ coordinates and let
$\mathcal O_\sigma=\{x_{\sigma(1)}\ge\cdots\ge
x_{\sigma(D)}\ge0\}$.  On this cone every $f_j$ is linear,
because the relative order of all coordinates within each block is
fixed.  Consequently $A\cap\mathcal O_\sigma$ is convex: it is
the intersection of a convex cone with the inverse image of $C$
under a linear map.

Let $S_k=\{\sigma(1),\ldots,\sigma(k)\}$ and let $i_j(k)$ be
the number of these coordinates in block $j$.  For the indicator
vector $\mathbf1_{S_k}$,
\[
 f_j(\mathbf1_{S_k}\vert_{E_j})=a_{j,i_j(k)}.
\]
The definition of $c_k$ thus implies
$v_k=\mathbf1_{S_k}/c_k\in A\cap\mathcal O_\sigma$.
It follows by convexity within this cone that
\[
 \Delta_\sigma=\operatorname{conv}(0,v_1,\ldots,v_D)
       \subset A\cap\mathcal O_\sigma.
\]
The columns $\mathbf1_{S_k}$ form a triangular matrix of
determinant $\pm1$ after reordering the coordinates.  Thus
\[
 \operatorname{vol}_D(\Delta_\sigma)
       =\frac1{D!\prod_{k=1}^D c_k}.
\]
The interiors of the $D!$ simplices are disjoint, because they lie
in the interiors of distinct coordinate-order cones.  Adding their
volumes gives $\operatorname{vol}_D(A)\ge1/\prod_kc_k$.
Together with \eqref{c6f:eq:A-volume}, this is exactly
\eqref{c6f:eq:packing-bound}.
\end{proof}

\subsection{The covering-product closure theorem}

\begin{lemma}[Block quotient lower bound]\label{c6f:lem:quotient}
Let $\Lambda_j$ be a full-rank lattice in $E_j$, and put
$\Lambda=\bigoplus_j\Lambda_j$.  Set
$a_{j,i}=\mu_i(K_j,\Lambda_j)$ and $a_{j,0}=0$.
Then the body \eqref{c6f:eq:aggregation} satisfies
\begin{equation}\label{c6f:eq:mu-lower}
 \mu_k(K_C,\Lambda)\ge c_k,
\end{equation}
where $c_k$ is given by \eqref{c6f:eq:convolution}.
\end{lemma}

\begin{proof}
For each $j$ choose a rational quotient $A_j:E_j\to E'_j$ of
rank $i_j$.  A rank-zero quotient is interpreted as a point with
gauge and covering radius zero.  Write $L_j=A_jK_j$ and
$\Gamma_j=A_j\Lambda_j$ in positive rank.  The quotient gauge is
\[
 g_{L_j}(u)=\min_{A_jx=u}g_j(x).
\]
The minimum is attained, since $g_j$ has compact sublevel sets.
Downward closure of $C$ consequently shows that the block quotient
$A=\bigoplus_jA_j$ satisfies $A(K_C)=(L_1,\ldots,L_m)_C$.

For a point $u=(u_j)$, minimize its gauge distance to
$\bigoplus_j\Gamma_j$.  Since $\gamma_C$ is coordinatewise
nondecreasing and each factor-lattice minimum is attained, the
minimum is
\[
 \gamma_C\left(
     \min_{z_1\in\Gamma_1}g_{L_1}(u_1-z_1),\ldots,
     \min_{z_m\in\Gamma_m}g_{L_m}(u_m-z_m)\right).
\]
Taking the supremum over the product of the factor tori gives the
exact covering-radius identity
\begin{equation}\label{c6f:eq:radius}
 \mu((L_1,\ldots,L_m)_C,\bigoplus_j\Gamma_j)
   =\gamma_C(\mu(L_1,\Gamma_1),\ldots,
                  \mu(L_m,\Gamma_m)).
\end{equation}
One may use maximizing points on the compact factor tori, or use
independent approximating sequences and continuity of $\gamma_C$.

Covering minima are the suprema of covering radii over rational
quotients of the appropriate rank.  Taking quotients $A_j$ whose
covering radii approach $a_{j,i_j}$ therefore proves
$\mu_k(K_C,\Lambda)\ge\gamma_C(a_{1,i_1},\ldots,a_{m,i_m})$
whenever $\sum_ji_j=k$.  Maximizing over these indices proves
\eqref{c6f:eq:mu-lower}.
\end{proof}

\begin{theorem}[Arbitrary monotone gauge aggregation]
\label{c6f:thm:aggregation}
Under the preceding hypotheses,
\begin{equation}\label{c6f:eq:supermultiplicativity}
 \boxed{\quad
 \mathcal F_D(K_C,\bigoplus_j\Lambda_j)
       \ge\prod_{j=1}^m\mathcal F_{d_j}(K_j,\Lambda_j).
 \quad}
\end{equation}
\end{theorem}

\begin{proof}
Divide \eqref{c6f:eq:body-volume} by
$\det\Lambda=\prod_j\det\Lambda_j$, and multiply by all the
covering minima.  Lemmas~\ref{c6f:lem:quotient}
and~\ref{c6f:lem:packing} give
\begin{align*}
 \mathcal F_D(K_C,\Lambda)
 &\ge W_{\boldsymbol d}(C)\prod_kc_k
       \prod_j\frac{\operatorname{vol}_{d_j}(K_j)}{\det\Lambda_j}\\
 &\ge\prod_j\left(
       \frac{\operatorname{vol}_{d_j}(K_j)}{\det\Lambda_j}
         \prod_{i=1}^{d_j}\mu_i(K_j,\Lambda_j)\right).
\end{align*}
No exact formula for the intermediate covering minima of $K_C$
is assumed or needed.
\end{proof}

\begin{corollary}[Closure of the conjectural class]
\label{c6f:cor:Fclosure}
If every factor satisfies
$\mathcal F_{d_j}(K_j,\Lambda_j)\ge(d_j+1)/2^{d_j}$, then
\[
 \mathcal F_D(K_C,\bigoplus_j\Lambda_j)
   \ge\frac{\prod_j(d_j+1)}{2^D}
   \ge\frac{D+1}{2^D}.
\]
The last inequality is strict when $m\ge2$.
\end{corollary}

\begin{proof}
Expand $\prod_j(1+d_j)$.  Its constant and linear terms sum to
$1+D$, and every higher-degree term is nonnegative.  For $m\ge2$
at least one quadratic term is positive.
\end{proof}

\begin{corollary}[Arbitrary aggregation of planar bodies and intervals]
\label{c6f:cor:planar-factors}
Suppose that each factor has dimension one or two, with $q$ planar
factors and any number of interval factors.  Then
\[
 \mathcal F_D(K_C,\bigoplus_j\Lambda_j)\ge(3/4)^q.
\]
In particular, for every planar convex body $K$, every interval
$I$ containing zero in its interior, and every two-dimensional
coefficient body $C$, the three-dimensional body
\[
 P=\{(x,t):(g_K(x),g_I(t))\in C\}
\]
satisfies $\mathcal F_3(P)\ge3/4>1/2$ for the product lattice.
\end{corollary}

\begin{proof}
An interval has covering product one, and Schnell's planar theorem
gives covering product at least $3/4$ for every planar convex body.
Apply Theorem~\ref{c6f:thm:aggregation}.  In arbitrary dimension,
$3^q2^{D-2q}\ge D+1$ also follows from
Corollary~\ref{c6f:cor:Fclosure}.
\end{proof}

For example, take any planar convex body containing zero in its
interior, an interval $[-a,b]$ with $a,b>0$, and a continuous concave
nonincreasing profile $\phi:[0,T]\to\mathbb R_{\ge0}$ with
$\phi(0)>0$.  Provided the region below the profile has nonempty
interior, the body
\[
 P=\{(x,t):0\le g_I(t)\le T,\quad
                         g_K(x)\le\phi(g_I(t))\}
\]
is covered by the corollary.  Its cross-sections are homothetic
copies of the arbitrary planar body $K$.  Nonlinear profiles give
families beyond Cartesian cylinders and free-sum bipyramids.

\subsection{A quantitative refinement for \texorpdfstring{$\ell_r$}{ell-r} gauge sums}

\begin{proposition}\label{c6f:prop:lp}
Let $p,q\ge1$, and let $K\subset\mathbb R^p$ and
$L\subset\mathbb R^q$ contain zero in their interiors.  For
$1\le r<\infty$ put
\[
 K\oplus_rL=\{(x,y):g_K(x)^r+g_L(y)^r\le1\}.
\]
For product lattices,
\begin{equation}\label{c6f:eq:lp-refinement}
 \mathcal F_{p+q}(K\oplus_rL)
 \ge C_r(p,q)\mathcal F_p(K)\mathcal F_q(L),
\end{equation}
where
\begin{equation}\label{c6f:eq:gamma-constant}
 C_r(p,q)=
 \frac{\Gamma(1+p/r)\Gamma(1+q/r)}{\Gamma(1+(p+q)/r)}
         \binom{p+q}{p}^{1/r}.
\end{equation}
The constant equals one at $r=1$, exceeds one for $1<r<\infty$,
and tends to one as $r\to\infty$.  At $r=\infty$ the body is
$K\times L$ and \eqref{c6f:eq:lp-refinement} holds with
$C_\infty(p,q)=1$.
\end{proposition}

\begin{proof}
Let $a_i=\mu_i(K)$ and $b_j=\mu_j(L)$, with $a_0=b_0=0$.
Lemma~\ref{c6f:lem:quotient} gives
\[
 \mu_k(K\oplus_rL)^r\ge
       \max_{i+j=k}(a_i^r+b_j^r).
\]
For positive sequences $u_i,v_j$ with $u_0=v_0=0$, their
maximum-sum convolution $w_k=\max_{i+j=k}(u_i+v_j)$ obeys
\begin{equation}\label{c6f:eq:maxplus}
 \prod_{k=1}^{p+q}w_k\ge\binom{p+q}{p}
                           \prod_{i=1}^p u_i\prod_{j=1}^q v_j.
\end{equation}
For completeness, induct on $N=p+q$, allowing empty sequences at
the boundary.  Deleting either last input bounds
$\prod_{k<N}w_k$ below by
\[
 \left(\prod_i u_i\prod_jv_j\right)
 \max\left\{\frac{\binom{N-1}{p-1}}{u_p},
             \frac{\binom{N-1}{p}}{v_q}\right\}.
\]
Multiply by $w_N=u_p+v_q$ and use
$\max(p/u_p,q/v_q)\ge N/(u_p+v_q)$.
This proves \eqref{c6f:eq:maxplus}.
Apply it with $u_i=a_i^r$ and $v_j=b_j^r$ and then take the
$r$th root.  The volume factor in \eqref{c6f:eq:body-volume} is
\[
 W_{(p,q)}(\{s,t\ge0:s^r+t^r\le1\})
  =\frac{\Gamma(1+p/r)\Gamma(1+q/r)}
          {\Gamma(1+(p+q)/r)},
\]
as follows by $u=s^r,v=t^r$ and the beta integral.  This proves
\eqref{c6f:eq:lp-refinement} with the stated constant.

Here is a complete check of its sign.  Put
\[
 h(t)=\log\Gamma(1+t(p+q))
        -\log\Gamma(1+tp)-\log\Gamma(1+tq),\quad t\ge0.
\]
We claim that $h$ is strictly convex.  For $x>0$, the trigamma
integral implies
\[
 x\psi'(1+x)=\int_0^\infty xe^{-xs}\frac{s}{e^s-1}\,ds.
\]
The function $s/(e^s-1)$ is strictly decreasing.  If $E$ is an
exponential random variable of mean one, the integral is
$\mathbb E[(E/x)/(e^{E/x}-1)]$ and is therefore strictly
increasing in $x$.  Hence
$G(x)=x^2\psi'(1+x)$ is strictly superadditive on positive
arguments: $G(x+y)>G(x)+G(y)$.  Consequently, for $t>0$,
\[
 h''(t)=\frac{G(t(p+q))-G(tp)-G(tq)}{t^2}>0.
\]
Since $h(0)=0$, strict convexity gives $h(t)<th(1)$ for
$0<t<1$.  Now
\[
 \log C_r(p,q)=r^{-1}h(1)-h(r^{-1}),
\]
which proves the strict assertion.  Both endpoint assertions
follow from the explicit gamma formula.  For $r=\infty$, use
Theorem~\ref{c6f:thm:aggregation} with $C=[0,1]^2$.
\end{proof}

In the three-dimensional quadratic case the numerical constant is
particularly simple.  For every planar $K$ and interval $I$ as above,
\[
 \mathcal F_3(K\oplus_2 I)
     \ge\frac{2\sqrt3}{3}\,\frac34
     =\frac{\sqrt3}{2}>\frac12.
\]
This is a lower bound, not a claimed optimal constant for that family.

\subsection{Exact spectrum preservation by asymmetric orthant changes}

\begin{proposition}\label{c6f:prop:asymmetric}
Let $C\subset\mathbb R_{\ge0}^d$ be any coefficient body and let
$a_i,b_i>0$.  Define
\[
 \eta_i(x)=\begin{cases}x/b_i,&x\ge0,\\-x/a_i,&x\le0,
                    \end{cases}
 \qquad
 K=\{x:(\eta_1(x_1),\ldots,\eta_d(x_d))\in C\}.
\]
Put $s_i=a_i+b_i$ and
\[
 U=\{x:(2|x_1|/s_1,\ldots,2|x_d|/s_d)\in C\}.
\]
With respect to $\mathbb Z^d$,
\begin{equation}\label{c6f:eq:spectrum}
 \operatorname{vol}_d(K)=\operatorname{vol}_d(U),\qquad
 \mu_k(K)=\mu_k(U)=
    \max_{|S|=k}\gamma_C\left(\sum_{i\in S}s_i^{-1}e_i\right)
       \quad(1\le k\le d).
\end{equation}
In particular, $\mathcal F_d(K)\ge1$.
\end{proposition}

\begin{proof}
The volume in an orthant is $\operatorname{vol}_d(C)$ times
the product of the corresponding $a_i$ or $b_i$.  Summing over
orthants proves
$\operatorname{vol}_d(K)=\operatorname{vol}_d(C)\prod_i s_i$;
the same formula holds for $U$.

Both bodies have the property that the projection onto any
coordinate subspace agrees with the intersection with that
subspace.  The reason is precisely the downward closure of $C$:
zeroing a coordinate keeps a point in the body.  Let $L_S$ be the
coordinate subspace supported on $S$.  Its intersection body is
the gauge aggregation of the intervals $[-a_i,b_i]$, $i\in S$,
with the remaining gauge coordinates zero.  Each of these intervals
has covering radius $1/s_i$.  The exact radius identity
\eqref{c6f:eq:radius} therefore gives
\[
 \mu(K\cap L_S,\mathbb Z^d\cap L_S)
       =\gamma_C\left(\sum_{i\in S}s_i^{-1}e_i\right).
\]
The coordinate projection description of covering minima gives
the lower inequality in \eqref{c6f:eq:spectrum}.

For the reverse inequality, put $t$ equal to the maximum of these
section radii over $|S|=k$.  Then $tK+\mathbb Z^d$ contains every
coordinate $k$-dimensional linear subspace: its subset
$t(K\cap L_S)+(\mathbb Z^d\cap L_S)$ covers $L_S$.
Every affine subspace $H$ of codimension $k$ meets at least one
such $L_S$.  To see this, write $H=\{x:Bx=c\}$ with $B$ a
rank-$k$ matrix and choose $k$ linearly independent columns of
$B$.  The corresponding $L_S$ maps onto $\mathbb R^k$, so
$H\cap L_S\ne\emptyset$.  Thus $tK+\mathbb Z^d$ meets
every such $H$, proving $\mu_k(K)\le t$.  This handles all
affine subspaces, including irrational ones.  The same argument
applies to $U$ and proves the exact spectrum identity.

Finally, apply Theorem~\ref{c6f:thm:aggregation} with all factors
one-dimensional intervals, whose covering products equal one.
Equivalently, the volume and spectrum identities transfer the
known standard-unconditional bound $\mathcal F_d(U)\ge1$.
\end{proof}

\subsection{Provenance and the remaining bridge}

The geometric proofs above are self-contained apart from the
standard rational-quotient characterization of covering minima;
the planar application also uses Schnell's theorem.  The exact
coordinate-section argument in the last proposition is the
mechanism in Krivoku\'ca's Theorem~1.3 and Corollary~3.3
(arXiv:2601.15173v1).  The standard-unconditional bound itself
is already Theorem~4.5 of Gonz\'alez Merino--Schymura, and is
not a newly solved case here.

The corpus contributed proof-design mechanisms, not unproved
substitutes for the geometric steps.  In manuscript \#185,
the theorem labelled \texttt{density-thm:weighted}, ``Exact
weighted optimum,'' proves a global inequality by assigning
nonnegative masses to lattice corners, verifying local capacities,
and summing.  Lemma~\ref{c6f:lem:packing} transfers this
certificate architecture to disjoint permutation simplices;
all required containment and determinant checks are given above.
In manuscript \#123, the section labelled \texttt{sec:hqmajor},
especially \texttt{hq:sharp} and \texttt{hq:finite}, uses
beta-integral slices and multiplicative gamma constants to
transfer inequalities between products.  The $\ell_r$ refinement
uses the same analytic design, with its own explicitly proved
gamma-constant inequality.  Neither corpus theorem is assumed
to imply a covering inequality merely by analogy.

The result does not represent an arbitrary convex body as
\eqref{c6f:eq:aggregation}.  Such a representation requires one
fixed gauge in each block and a common monotone coefficient body;
general fibers can change their shapes and centers.  Moreover,
one cannot infer a factor inequality from a proved inequality for
an aggregate: \eqref{c6f:eq:supermultiplicativity} points in the
opposite direction.  A solution of the unrestricted conjecture
would require a new reduction controlling both volume and the
entire covering-minimum product under replacement by an aggregate.
No such reduction is established here.  Historical priority for
the aggregation theorem or the asymmetric extension is not
asserted.

There is also a concrete obstruction to a volume-and-spectrum
preserving reduction to one-dimensional factors.  For the standard
terminal simplex $T_D$ the values $\mu_i(T_D)=i/2$ and
$\operatorname{vol}_D(T_D)=(D+1)/D!$ give
$\mathcal F_D(T_D)=(D+1)/2^D<1$ when $D\ge2$, whereas every
aggregation of intervals has covering product at least one.
More generally, a nontrivial aggregation of factors satisfying the
conjectured bounds has covering product strictly above the proposed
extremal value, by Corollary~\ref{c6f:cor:Fclosure}.  Thus the
theorem genuinely excludes these decompositions for an extremizer;
it does not make all bodies admit such a decomposition.


\section{Recalled input: simplex congruence states}
\label{eq:section}

We reproduce the ambient-lattice simplex criterion from the preceding
note, including its proof, because its actual states are the inputs
to the new polytope recursion in the next section.

Let $S=\operatorname{conv}(v_0,\ldots,v_d)$ be a lattice simplex with
$0\in\operatorname{int}S$. Put
\[
 \Gamma=\sum_{j=0}^{d}\mathbb Z v_j\subseteq\Lambda,
 \qquad \sum_{j=0}^{d}q_jv_j=0,
 \qquad \gcd(q_0,\ldots,q_d)=1,
 \qquad Q=\sum_{j=0}^{d}q_j,
\]
where every $q_j$ is positive.  If each $q_j$ divides $Q$, set
$m_j=Q/q_j$. Define the facet forms by
\begin{equation}\label{eq:facetforms}
 \ell_j(v_i)=1-m_j\delta_{ij}.
\end{equation}
The positive relation makes these well-defined linear forms, and
$S=\{x:\ell_j(x)\le1\text{ for every }j\}$.

\begin{lemma}[Deep points of the coarse lattice]\label{eq:deep}
For $x=\sum_jt_jv_j$ one has
\begin{equation}\label{eq:distance}
 \rho_{S,\Gamma}(x)=\min_{s\in[0,1)}
                       \sum_{j=0}^{d}\{t_j+s q_j\}.
\end{equation}
The maximum of this function is at most $d/2$. It equals $d/2$
at $x$ if and only if the coordinate crossing events are simple and
form an equally spaced grid of $Q$ points on the circle.
Consequently a point at distance $d/2$ exists if and only if every $m_j$ is
an integer and there are residues $0\le a_j<m_j$ whose classes
$a_j+m_j\mathbb Z$ partition $\mathbb Z$. Its class modulo $\Gamma$
then has a representative
\begin{equation}\label{eq:deep-representative}
 x(a)=-\sum_j\frac{a_j}{m_j}v_j,
 \qquad \ell_j(x(a))=a_j-\frac d2.
\end{equation}
In particular, every point at distance $d/2$ satisfies
$\ell_j(x)\equiv-d/2\pmod{\mathbb Z}$ for every $j$.
\end{lemma}
\begin{proof}
Every translate by $\Gamma$ is obtained by replacing $t$ by an
integer translate, and every other coefficient representation has
the form $t+s q$. Minimizing the nonnegative coefficient sum gives
\eqref{eq:distance}. Write $F(s)=\sum_j\{t_j+s q_j\}$.
It has slope $Q$, mean $(d+1)/2$, and downward jumps $J_e$ at
the distinct events $e$. The jump identity, obtained by integrating
$d(F^2)/ds$ on the event intervals, is
\[
 \frac1Q\sum_e J_eF(e)
 =\frac{d+1}{2}-\frac1{2Q}\sum_eJ_e^2.
\]
Since $\sum_eJ_e=Q$ and every $J_e$ is a positive integer, the
right side is at most $d/2$. Equality of the minimum with $d/2$
forces every jump to be one and every post-jump value to be $d/2$.
The slope then forces gaps $1/Q$. The converse follows from the
same identity. A coordinate has crossings at spacing $1/q_j$,
so its grid positions are one residue class modulo $m_j=Q/q_j$.
They partition the grid exactly when they partition $\mathbb Z$.

For these classes, take $t_j=-a_j/m_j$. The resulting events
are exactly that grid. Exactly one residue $a_j$ is zero, so
$F(0)=d-\sum_j a_j/m_j=d/2$. Thus
$\sum_j a_j/m_j=d/2$, and applying \eqref{eq:facetforms} to
$x(a)$ gives \eqref{eq:deep-representative}. Conversely an arbitrary
point at distance $d/2$ has such coefficients after a common rotation of the
event parameter and changes by integer coefficients; the former
adds a multiple of $q$ and does not change the represented point.
The facet forms are integral on $\Gamma$, giving the last assertion.
\end{proof}

Choose any lattice basis $b_1,\ldots,b_d$ of $\Lambda$. When the
forms $\ell_j$ are integral on $\Lambda$, define
\begin{equation}\label{eq:Dij}
 D_{ij}=\gcd\bigl(m_i,m_j,
       (\ell_i-\ell_j)(b_1),\ldots,(\ell_i-\ell_j)(b_d)\bigr)
       \qquad(0\le i<j\le d).
\end{equation}
The gcd is positive and independent of the chosen lattice basis.

\begin{theorem}[Ambient-lattice congruence criterion]\label{eq:main}
The equality $\mu(S,\Lambda)=d/2$ holds if and only if all three
conditions below hold:
\begin{enumerate}
\item every $q_j$ divides $Q$;
\item every $\ell_j$ is integral on $\Lambda$;
\item there are integers $a_0,\ldots,a_d$, equivalently residues
      $0\le a_j<m_j$, such that
\begin{equation}\label{eq:avoid}
       a_i\not\equiv a_j\pmod{D_{ij}}
       \quad\text{for every }i<j.
\end{equation}
\end{enumerate}
If these conditions hold, \eqref{eq:deep-representative} gives a
deepest point for the full ambient lattice. In particular, a pair
with $D_{ij}=1$ rules out extremality.
\end{theorem}
\begin{proof}
Suppose first that $\mu(S,\Lambda)=d/2$. The distance is continuous
and lattice-periodic, so it has a maximizing point $x$. Since
$\Gamma\subseteq\Lambda$, for every $h\in\Lambda$ one has
\[
 d/2\le\rho_{S,\Gamma}(x-h)\le d/2.
\]
Lemma~\ref{eq:deep} gives condition (1). It also gives
$\ell_j(x),\ell_j(x-h)\equiv-d/2\pmod{\mathbb Z}$, hence
$\ell_j(h)\in\mathbb Z$, proving (2).
Replace $x$ by a representative $x(a)$ modulo $\Gamma$ as in
the lemma; this does not change its property of being deepest
for $\Lambda$.

We next describe the events for $x(a)-h$. For a coefficient
representation $y=\sum_jt_jv_j$, let $A=\sum_jt_j$. Then
$\ell_j(y)=A-m_jt_j$. In particular, if
$\ell_j(y)=r_j-d/2$ with all $r_j$ integral, the crossings of
$\{t_j+s q_j\}$ satisfy
\[
 Qs+A+d/2\equiv r_j\pmod{m_j}.
\]
Thus its distance equals $d/2$ if and only if the residue classes
$r_j+m_j\mathbb Z$ form an exact cover. Applied to $y=x(a)-h$,
these are $r_j=a_j-\ell_j(h)$.

Two residue classes with moduli $m_i,m_j$ intersect exactly when
their residues agree modulo $g_{ij}=\gcd(m_i,m_j)$. Their density
sum over all $j$ is $\sum_j1/m_j=1$. Since all the moduli divide
$Q$, pairwise disjointness is therefore equivalent to exact
coverage, by counting the residues in one period of length $Q$.
Consequently the translated systems are exact covers for every
$h\in\Lambda$ precisely when
\[
 a_i-a_j\notin
 \{(\ell_i-\ell_j)(h)\bmod g_{ij}:h\in\Lambda\}
 \quad\text{for every }i<j.
\]
The subgroup on the right is the set of multiples of $D_{ij}$
modulo $g_{ij}$, by the gcd definition and Bezout's identity.
This proves (3).

Conversely assume (1)--(3). Since $D_{ij}\mid g_{ij}$,
\eqref{eq:avoid} first implies that $a_j+m_j\mathbb Z$ is an exact
cover. Choose $x=x(a)$. The same subgroup argument shows that
$a_j-\ell_j(h)+m_j\mathbb Z$ remains an exact cover for every
$h\in\Lambda$. By the event calculation,
$\rho_{S,\Gamma}(x-h)=d/2$ for every such $h$. Finally
\[
 \rho_{S,\Lambda}(x)
 =\min_{h\in\Lambda/\Gamma}\rho_{S,\Gamma}(x-h)=d/2.
\]
The universal upper bound proves the desired equality.
\end{proof}

\begin{remark}[Relation to the earlier criterion]\label{eq:relation}
If $\Lambda=\Gamma$, the numbers in \eqref{eq:Dij} coming from
vertex generators are $0$, $-m_i$, or $m_j$. Hence
$D_{ij}=\gcd(m_i,m_j)$, and the theorem reduces to the exact
covering criterion of \cite{VMRC200}. For a strict superlattice
the smaller gcd values encode precisely the additional lattice
translates that may destroy a former deepest point. The theorem
is finite: after integral linear algebra one tests at most
$\prod_jm_j$ residue tuples. It makes no efficiency or polynomial
time claim.
\end{remark}



\section{Every extremal polytope has finitely many deep holes}
\label{ce:section}

For a convex body $K$ containing the origin in its interior, write
$g_K$ for its gauge and put
\[
 \rho_{K,\Lambda}(x)=\min_{z\in\Lambda}g_K(x-z),\qquad
 \mathcal H(K,\Lambda)=
 \{x+\Lambda:\rho_{K,\Lambda}(x)=d/2\}.
\]
Here $d$ is the dimension of $K$. Thus $\mathcal H$ is empty unless
the non-hollow upper bound is attained, when $K$ is a non-hollow
lattice polytope. We count classes in the lattice torus, not all
their lattice translates.
Rationality below means rational coordinates in a basis of $\Lambda$.

The conditional-state composition in corpus [041] suggests retaining
the actual finite boundary states through a geometric projection,
rather than retaining only a scalar covering bound. The states here
are exact-covering-system representatives of simplex deep holes.
All assertions about their transfer are proved below; no Ewald
counting theorem is an input.

\begin{lemma}[Counting the simplex states]\label{ce:simplex-count}
Let $S$ be a lattice $r$-simplex containing $0$ in its interior.
Let $q_0,\ldots,q_r$ be its primitive positive vertex relation,
$Q=\sum q_j$, and $\Gamma$ its vertex-generated lattice.
If $\mu(S,\Lambda)=r/2$, then
\[
 |\mathcal H(S,\Lambda)|
 \le \frac{(r+1)!}{Q[\Lambda:\Gamma]}\le r!.
\]
Equality in the last bound is possible only for the standard
terminal simplex in its vertex-generated lattice, and it holds
for that simplex.
\end{lemma}
\begin{proof}
The simplex congruence criterion gives integers $m_j=Q/q_j$.
In the coarse lattice, its deep-hole representatives correspond
to exact covers by the labelled residue classes
$a_j+m_j\mathbb Z$, with $0\le a_j<m_j$.
There are at most $(r+1)!$ such labelled covers. Indeed, the
least nonnegative representatives $a_j$ are all distinct. Their
increasing order determines the cover uniquely: its first residue
is zero, and, after a collection of classes has been assigned,
the next residue is the least nonnegative integer not covered by
those classes. If that integer is not smaller than the modulus
of the next prescribed label, the prescribed order is impossible.
For a genuine exact cover, this reconstruction succeeds at every
step. Hence assigning the order is injective.

Simultaneous translation $a_j\mapsto a_j+k\pmod{m_j}$ acts
freely on these covers with period $Q$, since
$\operatorname{lcm}(m_0,\ldots,m_r)=Q$ follows from
$\gcd(q_0,\ldots,q_r)=1$. It does not change the represented
point modulo $\Gamma$: the change before reducing residues is
$-k\sum_j v_j/m_j=0$. Conversely, if $x(a)-x(b)\in\Gamma$,
then $(b_j-a_j)/m_j=n_j+tq_j$ for integers $n_j$ and one real
$t$, because the kernel of the vertex map is spanned by $q$.
Multiplication by $m_j$ gives
$b_j-a_j=n_jm_j+tQ$. Thus $k=tQ$ is an integer, and the
two residue arrays differ by a common integer shift.
Consequently there are at most $(r+1)!/Q$ coarse deep-hole classes.

A point is deep for $\Lambda$ precisely when all its translates
by $\Lambda/\Gamma$ are coarse deep holes. The surviving coarse
classes form free orbits of size $[\Lambda:\Gamma]$. This gives
the first inequality. Since $Q\ge r+1$, the second follows.
If the final bound is attained, $Q=r+1$ and
$[\Lambda:\Gamma]=1$, so every $q_j=1$. The simplex is then
unimodularly equivalent to
$T_r=\operatorname{conv}(-\mathbf1_r,e_1,\ldots,e_r)$.
For $T_r$ all $m_j=r+1$; all $(r+1)!$ permutations of the residues
form exact covers, and their common-shift orbits give exactly $r!$
deep-hole classes.
\end{proof}

\begin{lemma}[Finite-state lifting]\label{ce:lifting}
Let $P$ be a non-hollow lattice $d$-polytope with $0$ interior.
Let $S$ be a simplex on vertices of $P$ with
$0\in\operatorname{relint}S$, put $r=\dim S<d$, and let
$L=\operatorname{span}S$, $\Lambda_L=\Lambda\cap L$,
$\pi:\mathbb R^d\to\mathbb R^d/L$, $Q=\pi P$.
If $\mu(P,\Lambda)=d/2$, then
\[
 P\cap L=S,\quad \mu(S,\Lambda_L)=r/2,\quad
 \mu(Q,\pi\Lambda)=(d-r)/2.
\]
Moreover, if the two lower-dimensional deep-hole sets are finite,
\[
 |\mathcal H(P,\Lambda)|
 \le |\mathcal H(S,\Lambda_L)|\,
       |\mathcal H(Q,\pi\Lambda)|.
\]
\end{lemma}
\begin{proof}
Set $\alpha=r/2$ and $\beta=(d-r)/2$. The section--projection
bound and the non-hollow upper bound give
\[
 d/2=\mu(P)\le\mu(P\cap L)+\mu(Q)
 \le\mu(S)+\mu(Q)\le\alpha+\beta=d/2.
\]
Every inequality is equality. Tightness of a simplex, also for
nonlattice convex super-bodies, gives $P\cap L=S$.

Fix a representative $y$ of each class in $\mathcal H(Q,\pi\Lambda)$.
Choose $h_y\in\pi\Lambda$ and $v_y\in\beta P$ such that
$\pi v_y=y-h_y$, and choose a lattice lift $z_y$ of $h_y$.
Put $w_y=v_y+z_y$, so $\pi w_y=y$. Such choices exist because
$\beta Q+\pi\Lambda$ covers the quotient.

If $x$ is deep for $P$, then $\pi x$ is deep for $Q$. Otherwise,
lifting a covering point for $\pi x$ at scale strictly smaller
than $\beta$ and completing along $S$ at scale $\alpha$ would
cover $x$ at scale strictly smaller than $d/2$.
After a lattice translation, arrange $\pi x=y$. Since
$x-w_y\in L$, convexity and $S\subseteq P$ imply
\[
 \rho_{P,\Lambda}(x)
 \le\beta+\rho_{S,\Lambda_L}(x-w_y)
 \le\beta+\alpha.
\]
Equality of the outer terms forces $x-w_y$ to be deep for $S$.
Thus every class of $\mathcal H(P,\Lambda)$ has a representative
$w_y+s$, with $s$ chosen from the finite set of simplex states.
There are at most the displayed number of candidates.
\end{proof}

\begin{theorem}[A sharp factorial bound and its rigidity]
\label{ce:factorial}
For every non-hollow lattice $d$-polytope $P$,
\[
 |\mathcal H(P,\Lambda)|\le d!.
\]
All points in this set are rational modulo $\Lambda$ and isolated.
If $P$ is not a simplex and $d\ge2$, the stronger bound
\[
 |\mathcal H(P,\Lambda)|\le(d-1)!
\]
holds, and this stronger bound is also sharp.
Equality $|\mathcal H(P,\Lambda)|=d!$ holds if and only if
$P$, after translating its interior lattice point to the origin,
is unimodularly equivalent to $T_d$.
\end{theorem}
\begin{proof}
The assertion is empty if $\mu(P)<d/2$. Otherwise, choose an
inclusion-minimal set of vertices whose convex hull contains zero.
It is a simplex $S$ with zero in its relative interior.
If $\dim S=d$, tightness and $\mu(S)\le d/2$ force $P=S$;
Lemma~\ref{ce:simplex-count} applies.
If $r=\dim S<d$, apply Lemma~\ref{ce:lifting} and induction:
\[
 |\mathcal H(P,\Lambda)|\le r!(d-r)!\le(d-1)!<d!.
\]
The penultimate inequality follows from $\binom dr\ge d$.
The simplex representatives are rational by their exact-cover
formula. In the inductive construction, $y$ is rational and a
nonempty rational polytope fiber admits a rational $v_y$.
Thus all candidates, and hence all surviving deep holes, are
rational. A finite subset of the torus is isolated. The rigidity
statement now follows from Lemma~\ref{ce:simplex-count}.
For sharpness of the nonsimplicial bound, take
$T_{d-1}\oplus[-1,1]$ in the product lattice. Its gauge and
lattice distance are sums of the two factor gauges and distances,
so its deep-hole set is their Cartesian product. It has exactly
$(d-1)!$ deep-hole classes and is not a simplex.
\end{proof}

For comparison, the extremal four-dimensional simplex with relation
$(1,1,1,1,2)$ has exactly $12$ deep-hole classes, not $4!=24$.
Its moduli are $(6,6,6,6,3)$: the modulus-three class has three
choices, and the four remaining classes occupy the four unclaimed
residues in $4!$ orders. Thus there are $3\cdot4!=72$ labelled
exact covers and $72/6=12$ common-shift orbits. This simplex
attains the same covering radius as $T_4$ while the number of
last-covered classes distinguishes them.

\subsection{An exact decision procedure for arbitrary polytopes}

The preceding proof yields a finite recognition criterion for
\emph{all} non-hollow lattice polytopes, including nonsimplicial
ones. It also lists every deep-hole class when equality holds.
It is not a short geometric classification of all extremizers.

\begin{theorem}[Recursive equality certificate]\label{ce:algorithm}
The following rational finite procedure returns exactly
$\mathcal H(P,\Lambda)$.
\begin{enumerate}
\item Translate an interior lattice point to zero, and choose a
      minimal positive vertex simplex $S$.
\item If $S$ is full dimensional, return the simplex congruence
      states when $P=S$, and the empty set when $P\ne S$.
\item Otherwise, form $L$, $Q$, and the two induced lattices.
      Return the empty set if $P\cap L\ne S$, if $S$ has no
      extremal congruence state, or if the recursive call on $Q$
      returns the empty set.
\item Choose one rational lift $w_y\in\beta P+\Lambda$ of each
      quotient state as in Lemma~\ref{ce:lifting}. Form all
      candidates $c=w_y+s$ using the simplex states $s$.
\item Retain $c$ exactly when
      \begin{equation}\label{ce:finite-empty}
       \operatorname{int}\bigl(c-(d/2)P\bigr)\cap\Lambda
       =\varnothing.
      \end{equation}
      Remove duplicates modulo $\Lambda$ and return the result.
\end{enumerate}
In particular, equality $\mu(P,\Lambda)=d/2$ is equivalent to
the returned set being nonempty.
\end{theorem}
\begin{proof}
All rejections are necessary by simplex tightness and
Lemma~\ref{ce:lifting}. That lemma proves that every deep hole
occurs among the candidates. Conversely,
\eqref{ce:finite-empty} says $\rho_{P,\Lambda}(c)\ge d/2$;
the universal upper bound makes it an equality.
Each recursion reduces dimension. The simplex tests use only
integral linear algebra and finite residue enumeration. Rational
polyhedral operations give the sections, quotients, and lifts.
Finally, the body in \eqref{ce:finite-empty} is bounded: enumerate
the finitely many lattice points in a rational coordinate bounding
box and test its strict facet inequalities exactly. This proves
termination and exactness, with no efficiency claim.
\end{proof}

\begin{remark}[A useful obstruction]
The same argument gives a geometric rejection test before the
last enumeration. If $P$ is extremal, $y$ is a quotient deep hole
which is the projection of a deep hole of $P$, and
$u\in\beta Q\cap(y-\pi\Lambda)$, then the fiber
$\beta P\cap\pi^{-1}(u)$ is a singleton. Indeed all its points
$v$ would make $x-z-v$ simplex deep holes, for a fixed deep lift
$x$ and fixed lattice lift $z$ of $y-u$. A positive-dimensional
convex fiber would give infinitely many points in the discrete
set of simplex deep holes, a contradiction. A quotient state
without a deep lift is not subject to this obstruction.
\end{remark}

\paragraph{Relation to existing algorithms.}
Finite computability of a rational polytope's covering radius is
already known; see Cslovjecsek--Malikiosis--Nasz\'odi--Schymura
\cite{ComputingRadius}, and the earlier algorithm of Kannan cited
there. The contribution proved here is the specialized recursive
state description at the non-hollow extremal value, the sharp
dimension-only count of all such states, and its rigidity statement.
No first algorithm or improvement in worst-case running time is
claimed.


\appendix
\section{Recalled input: the one-small-column planar theorem}
\label{app:onesmall}
This earlier lemma is included to make the three-weight proof
independent of access to the preceding manuscript.

\begin{lemma}[Three first-event bounds]
\label{r4w:lem-first-events}
Let $p,q\ge1$ be integers and $0<b\le1$. For all real $t_0,t_1,t_2$,
\[
 \min_{s\in\mathbb R}
 \bigl(\{t_0+s\}+b\{t_1+ps\}+b\{t_2+qs\}\bigr)
 \le \frac{b(2+b)}{1+2b}.
\]
\end{lemma}
\begin{proof}
Rotate the parameter to make $t_0=0$ modulo one, and replace the other
phases by their representatives in $[0,1)$. Their first subsequent
crossing times are
\[
 \alpha=\frac{1-t_1}{p},\qquad
 \beta=\frac{1-t_2}{q}.
\]
Interchanging the two equally weighted coordinates, if needed, makes
$\alpha\le\beta$. Both times belong to $(0,1]$. Put
$C=1+b(p+q)$ and denote the displayed fractional-part sum by $F$.
Then
\begin{align*}
 F(0)&=L_0:=2b-bp\alpha-bq\beta,\\
 F(\alpha)&\le L_1:=b+(1+bq)\alpha-bq\beta,\\
 F(\beta)&\le L_2:=(1+bp)\beta-bp\alpha.
\end{align*}
The second estimate subtracts at least the first $p$-coordinate jump.
The third subtracts at least one jump of each of the two coordinates.
Any extra crossings only strengthen these inequalities, including
coincidences and a crossing of the first coordinate at time one.

The positive numbers $1/C,bp/C,bq/C$ sum to one. Their corresponding
weighted average of $L_0,L_1,L_2$ is
\[
 \frac{L_0+bpL_1+bqL_2}{C}
 =\frac{b(2+bp)}{1+b(p+q)}.
\]
Thus the minimum of $F$ is at most this last quantity. Finally,
\[
 \frac{b(2+b)}{1+2b}
 -\frac{b(2+bp)}{1+b(p+q)}
 =\frac{b^2\bigl((1-b)(p-1)+(2+b)(q-1)\bigr)}
 {(1+2b)(1+b(p+q))}\ge0.
\]
This proves the lemma without a restriction on later crossings.
\end{proof}

\begin{lemma}[Planar unimodular refinement inside a triangle]
\label{r4w:lem-fan}
Let $v,w$ be independent vectors of a planar lattice $\Lambda$.
Their cone admits a subdivision into unimodular cones whose primitive
ray generators all lie in $\operatorname{conv}(0,v,w)$.
\end{lemma}
\begin{proof}
First shorten the two boundary vectors to their primitive generators.
Suppose a current pair of primitive generators $u_1,u_2$ has lattice
determinant $D>1$. A nonzero lattice class in their half-open
fundamental parallelogram has a representative
$z=\alpha u_1+\beta u_2$ with $0\le\alpha,\beta<1$.
Primitivity of the two generators forces $0<\alpha,\beta<1$.
If $\alpha+\beta\le1$, retain $z$; otherwise retain $u_1+u_2-z$.
The retained vector belongs to $\operatorname{conv}(0,u_1,u_2)$
and lies strictly between the boundary rays. Shorten it to its
primitive generator and subdivide along that ray.
Both new positive lattice determinants are strictly smaller than $D$.
Induction terminates in determinant-one cones. Every new triangle
lies in its predecessor, so all generators remain in the original
triangle.
\end{proof}

\begin{theorem}[Planar stretching with at most one fixed column]
\label{r4w:thm-stretch}
Let $u_1,\ldots,u_N$ be vectors of a planar lattice $\Lambda$,
with $0\in\operatorname{int}\operatorname{conv}(u_1,\ldots,u_N)$.
Keep at most one of the columns fixed and multiply all the others
by $r\ge1$. For their resulting convex hull $P_r$,
\[
 \boxed{\displaystyle
 \mu(P_r,\Lambda)\le G(r):=\frac{2r+1}{r(r+2)}.}
\]
The columns need not be primitive, distinct, or generators of $\Lambda$.
\end{theorem}
\begin{proof}
The case $r=1$ is included below. Delete zero columns, which does not
affect origin-interiority or the convex hull. A minimally supported
positive dependence among the columns consists either of an opposite
pair or of a triangle containing zero in its interior.

First consider such a triangle. If none of its vertices is fixed,
its scaled hull has covering radius at most $1/r\le G(r)$ by the
planar non-hollow bound. If exactly one is fixed, write its hull as
\[
 T_r=\operatorname{conv}(v_0,rv_1,rv_2),
 \qquad 0\in\operatorname{int}\operatorname{conv}(v_0,v_1,v_2).
\]
The vectors $v_1,v_2$ are independent, and $-v_0$ lies strictly
inside their positive cone. Apply Lemma~\ref{r4w:lem-fan} to that cone.

If $-v_0$ lies strictly inside one of its unimodular cones, let
$w_1,w_2$ be the generators of that cone. Then
\[
 v_0=-pw_1-qw_2\quad\text{for integers }p,q\ge1,
\]
and
$T'_r=\operatorname{conv}(v_0,rw_1,rw_2)\subseteq T_r$.
The inclusion holds because $rw_j\in\operatorname{conv}(0,rv_1,rv_2)$
and $0\in T_r$. For any point written in the integer directions
$v_0,w_1,w_2$, fractional-part rounding along their relation
$(1,p,q)$ bounds its distance to $\Lambda$ by
\[
 \min_s\bigl(\{t_0+s\}+r^{-1}\{t_1+ps\}
                    +r^{-1}\{t_2+qs\}\bigr).
\]
Lemma~\ref{r4w:lem-first-events}, with $b=1/r$, bounds this by
$G(r)$, as required.

If $-v_0$ lies on a ray of the subdivision, write $v_0=-Aw$
with $w$ primitive and $A\ge1$ an integer. The section of $T_r$
along $\mathbb Rw$ contains $[-Aw,rw]$, of lattice length at
least $r+1$. The quotient of $T_r$ by this line has lattice
length at least $2r$: its two nonzero endpoint candidates are
$r$ times opposite-sign nonzero integers. Section/projection
subadditivity therefore gives
\[
 \mu(T_r,\Lambda)\le\frac1{r+1}+\frac1{2r}\le G(r),
\]
where
\[
 G(r)-\left(\frac1{r+1}+\frac1{2r}\right)
 =\frac{r-1}{2(r+1)(r+2)}\ge0.
\]

It remains to handle an opposite pair in the original columns.
Let $W$ be its line. If one selected column is fixed, the section
of $P_r$ along $W$ has lattice length at least $r+1$, and all
nonzero projected columns are scaled, so the quotient interval
has lattice length at least $2r$. If both selected columns are
scaled, the section length is at least $2r$. At most one nonzero
projected column is fixed; because the unscaled quotient has
opposite-sign nonzero integer endpoints, the scaled quotient has
length at least $r+1$. Thus both cases give the same preceding
bound. Each selected triangle or section lies in $P_r$, so
monotonicity completes the proof.
\end{proof}



\section{Checks and the precise remaining questions}

\subsection{Mathematical and exact arithmetic checks}

The arguments are analytic proofs. The residue bounds include
coincident jumps, the sharp event schedules construct actual phases,
and the weighted upper bound ranges over every rational planar
quotient. The gauge-product argument proves containment separately
inside each coordinate-order cone; it does not assume that its
auxiliary piecewise-linear function is globally convex. The
deep-hole recursion uses the full ambient lattice at every stage.

The following finite checks were also performed with exact rational
arithmetic:
\begin{itemize}
\item The new arithmetic bound was checked on 2,500 unit-weight
flows and 2,500 commensurable weighted flows, including coincident
events. The sharp schedule for $(1^k,a^\ell)$ was checked on
all $1\le k,\ell,a\le10$, another 1,000 cases.
\item The two-event formula was checked on 1,008 exact optimization
cases and 6,552 actual rational phase certificates, followed by
243 integer-ratio cases and seven named examples.
\item Exact-cover enumeration gave $d!$ deep-hole classes for $T_d$,
$1\le d\le6$, and 12 classes for the four-dimensional relation
$(1,1,1,1,2)$, consistent with the direct count in the text.
\end{itemize}
These checks do not replace the general proofs. Independent readings
checked the event-deficit argument, the two-event attainment and
dimension truncation, the permutation-simplex packing, and the
deep-hole lifting argument. The proof text incorporates the
zero-column qualification in the planar reduction and the
deep-lift qualification in the fiber obstruction.

\subsection{What is still not established}

The new arithmetic theorem proves D whenever a suitable simplex
section satisfies $(h+1)G(q)\le2Q$. It does not prove that such a
certificate always exists. The relation $(1,3,4)$ already shows
why the gcd bound alone cannot supply an unrestricted proof.
The prime-volume theorem and repeated-relation exact family are
fully established without that further reduction.

The second covering minimum is completely determined here for the
specified three-weight family, including irrational weights.
The argument eliminating unequal planar frequencies uses $c\ge2b$.
It does not imply the missing interval $b<c<2b$, a formula for
independently varying later weights, or general higher covering
indices.

The gauge-aggregation theorem proves F for the stated geometric
class. An arbitrary convex body need not have homothetic sections
with fixed factor gauges and one common monotone coefficient body.
A reduction of general bodies to that class would need to control
both volume and all covering minima; it is not supplied here.
The obstruction at the end of the aggregation section explains why
simple spectrum-preserving reduction to intervals cannot work.

The recursive equality criterion applies to every non-hollow
lattice polytope and terminates. It decides equality and enumerates
all last-covered classes when equality occurs. A concise geometric
description of all extremal polytopes, and a general reflexivity
theorem for them, are not proved. General finite algorithms for
covering radius are existing results \cite{ComputingRadius}; the
specific state structure and sharp counts are what is proved here.

\subsection{Reproducibility and status}

The manuscript is one standalone \LaTeX{} source with its bibliography
included. All new estimates, sharpness constructions, and finite
reductions are given in the source. The only external geometric
inputs beyond definitions are standard covering-minimum quotient
facts and published planar covering-product bounds; the basic
flow, section argument, simplex criterion, and planar stretching
input are reproduced.

No universal literature-priority claim is made for these transfers.
The referenced papers fix the original questions and the established
ingredients. The scope labels refer to what has been proved in this
manuscript and the preceding work, rather than an assertion of external
peer review, formal verification, or resolution of every conjecture
in the source paper.

\begin{thebibliography}{99}
\bibitem{CSS}
G. Codenotti, F. Santos, and M. Schymura,
\emph{The covering radius and a discrete surface area for non-hollow
simplices}, Discrete Comput. Geom. 67 (2022), 65--111.
\href{https://doi.org/10.1007/s00454-021-00330-3}{doi:10.1007/s00454-021-00330-3}.
Open text: \href{https://pmc.ncbi.nlm.nih.gov/articles/PMC8709830/}{PMC8709830}.

\bibitem{GMS}
B. Gonz\'alez Merino and M. Schymura,
\emph{On densities of lattice arrangements intersecting every
$i$-dimensional affine subspace}, Discrete Comput. Geom. 58
(2017), 663--685.
\href{https://doi.org/10.1007/s00454-017-9911-x}{doi:10.1007/s00454-017-9911-x}.
The preprint lists the second author as Matthias Henze.

\bibitem{Schnell}
U. Schnell, \emph{A Minkowski-type theorem for covering minima
in the plane}, Geom. Dedicata 55 (1995), 247--255.
\href{https://doi.org/10.1007/BF01266316}{doi:10.1007/BF01266316}.

\bibitem{Krivokuca}
K. Krivoku\'ca, \emph{Upper bounds on covering minima of convex bodies},
arXiv:2601.15173v1 (21 January 2026).
\href{https://arxiv.org/html/2601.15173v1}{Full text}.

\bibitem{ComputingRadius}
J. Cslovjecsek, R. D. Malikiosis, M. Nasz\'odi, and M. Schymura,
\emph{Computing the covering radius of a polytope with an application
to lonely runners}, arXiv:2009.12080v2 (2021).
\href{https://real.mtak.hu/152905/1/2009.12080.pdf}{Primary manuscript}.


\bibitem{VMRC200}
Verifiable Mathematical Reasoning Corpus, manuscript [200], \emph{Weighted fractional-part flows and sharp lattice covering bounds}. 
\href{https://github.com/QEEDDD/Verifiable-Mathematical-Reasoning-Corpus/blob/ac1492a50784be919fe0b5e8f19f96f8d221a22e/\%5B200\%5D\%20lattice_covering_flows.tex}{Frozen source at revision \texttt{ac1492a5}}.

\bibitem{Corpus030}
Verifiable Mathematical Reasoning Corpus, manuscript [030], \emph{Arithmetic and observation classifications, recovery, and transcendence barriers}. 
\href{https://github.com/QEEDDD/Verifiable-Mathematical-Reasoning-Corpus/blob/ac1492a50784be919fe0b5e8f19f96f8d221a22e/\%5B030\%5D\%20Arithmetic\%20and\%20Observation\%20Classifications\%2C\%20Recovery\%2C\%20and\%20Transcendence\%20Barriers.tex}{Frozen source at revision \texttt{ac1492a5}}.

\bibitem{Corpus029}
Verifiable Mathematical Reasoning Corpus, manuscript [029], \emph{Corpus manuscript [029], simultaneous returns of finitely many phases}. 
\href{https://github.com/QEEDDD/Verifiable-Mathematical-Reasoning-Corpus/blob/ac1492a50784be919fe0b5e8f19f96f8d221a22e/\%5B029\%5D\%20Observation\%20and\%20Recovery\%20\%28Geometry\%2C\%20Mixtures\%2C\%20and\%20Spectral\%20Inversion\%29.tex}{Frozen source at revision \texttt{ac1492a5}}.

\bibitem{Corpus041}
Verifiable Mathematical Reasoning Corpus, manuscript [041], \emph{Ewald points of monotone polytopes and branching simplex bundles}. 
\href{https://github.com/QEEDDD/Verifiable-Mathematical-Reasoning-Corpus/blob/ac1492a50784be919fe0b5e8f19f96f8d221a22e/\%5B041\%5D\%20Ewald\%20points\%20of\%20monotone\%20polytopes\%20the\%20minimum\%20in\%20dimension\%20eight\%20and\%20branching\%20simplex\%20bundles.tex}{Frozen source at revision \texttt{ac1492a5}}.

\bibitem{Corpus123}
Verifiable Mathematical Reasoning Corpus, manuscript [123], \emph{Neighborhoods, kernels, and fractional zeros}. 
\href{https://github.com/QEEDDD/Verifiable-Mathematical-Reasoning-Corpus/blob/ac1492a50784be919fe0b5e8f19f96f8d221a22e/\%5B123\%5D\%20Neighborhoods_Kernels_and_Fractional_Zeros.tex}{Frozen source at revision \texttt{ac1492a5}}.

\bibitem{Corpus185}
Verifiable Mathematical Reasoning Corpus, manuscript [185], \emph{Three proof transfers unified}. 
\href{https://github.com/QEEDDD/Verifiable-Mathematical-Reasoning-Corpus/blob/ac1492a50784be919fe0b5e8f19f96f8d221a22e/\%5B185\%5D\%20three_proof_transfers_unified.tex}{Frozen source at revision \texttt{ac1492a5}}.

\bibitem{Corpus215}
Verifiable Mathematical Reasoning Corpus, manuscript [215], \emph{Two frontiers: spectral realization and tiling}. 
\href{https://github.com/QEEDDD/Verifiable-Mathematical-Reasoning-Corpus/blob/ac1492a50784be919fe0b5e8f19f96f8d221a22e/\%5B215\%5D\%20two_frontiers.tex}{Frozen source at revision \texttt{ac1492a5}}.

\end{thebibliography}
\end{document}
